% =========================================================================
% Problem Formulation + Proposed Architecture (GNN.py -> "CD-GNN")
% -------------------------------------------------------------------------
% GNN-version rewrite of formulation_and_architecture.tex for the RA-L
% submission built on SchedulerGNN (Static_alogorithm/GNN/GNN.py).
% Usage: % =========================================================================
% Problem Formulation + Proposed Architecture (GNN.py -> "CD-GNN")
% -------------------------------------------------------------------------
% GNN-version rewrite of formulation_and_architecture.tex for the RA-L
% submission built on SchedulerGNN (Static_alogorithm/GNN/GNN.py).
% Usage: % =========================================================================
% Problem Formulation + Proposed Architecture (GNN.py -> "CD-GNN")
% -------------------------------------------------------------------------
% GNN-version rewrite of formulation_and_architecture.tex for the RA-L
% submission built on SchedulerGNN (Static_alogorithm/GNN/GNN.py).
% Usage: % =========================================================================
% Problem Formulation + Proposed Architecture (GNN.py -> "CD-GNN")
% -------------------------------------------------------------------------
% GNN-version rewrite of formulation_and_architecture.tex for the RA-L
% submission built on SchedulerGNN (Static_alogorithm/GNN/GNN.py).
% Usage: \input{formulation_and_architecture_gnn} inside an IEEEtran doc.
% Requires in the preamble:
%   \usepackage{amsmath,amssymb,bm}
%   \usepackage{booktabs,array,multirow,graphicx}
%   \usepackage{algorithm}
%   \usepackage{algpseudocode}
% Citation keys are placeholders -- suggested BibTeX entries are listed in
% comments at the end. "CD-GNN" (Cross-Docking GNN) is a placeholder model
% name; rename globally if you prefer.
% Sections: Problem Formulation, Architecture, Surrogate Calibration,
% Policy Optimization (incl. post-processing local search).
% =========================================================================

\section{Problem Formulation}
\label{sec:problem}

\begin{table*}[t]
\caption{Summary of Notation}
\label{tab:notation}
\centering
\footnotesize
\begin{tabular}{@{}l p{0.295\textwidth} @{\qquad} l p{0.295\textwidth}@{}}
\toprule
Symbol & Meaning & Symbol & Meaning \\
\midrule
\multicolumn{2}{@{}l}{\itshape Problem setting} &
\multicolumn{2}{@{}l}{\itshape Architecture} \\
$\mathcal{W}$ & grid workspace (four-connected unit cells) &
  $\bm{x}^{A}_k,\bm{x}^{J}_j$ & AMR / job input features ($\mathbb{R}^{8},\mathbb{R}^{16}$) \\
$\Delta(u,v)$ & Manhattan distance between positions $u,v$ &
  $h$ & hidden width of all embeddings ($h{=}128$) \\
$\mathcal{J},\ n$ & job set; number of jobs &
  $\bm{h}^{A}_k$ & AMR embedding \\
$\mathcal{M},\ m$ & AMR fleet; fleet size &
  $\bm{g}^{(l)}_j$ & job embedding after $l$ GIN layers \\
$\mathcal{T}$ & set of material types &
  $\mathcal{G}_t=(\mathcal{J},E_t)$ & dynamic precedence graph at step $t$ \\
$\mathcal{D}^{\mathrm{in}},\mathcal{D}^{\mathrm{out}},\mathcal{D}$ & inbound docks; outbound docks; their union &
  $\mathcal{P}(j)$ & committed predecessors of job $j$ in $\mathcal{G}_t$ \\
$\tau_j$ & material type of job $j$ &
  $\epsilon^{(l)},\ L$ & learnable GIN self-weight; number of GIN layers \\
$o_j,\ d_j$ & inbound and outbound dock of job $j$ &
  $\mathrm{LN}$ & layer normalization \\
$r_j$ & release time of job $j$ &
  $\bm{e}_{\omega}$ & learned operation-type embedding \\
$p_j=p(\tau_j)$ & dock service duration of job $j$ &
  $z_{\omega,k,j}$ & logit of action $(\omega,k,j)$ \\
$\pi_j,\ \delta_j$ & pickup and delivery operations of job $j$ &
  $V_{\phi}$ & critic (state-value) network, parameters $\phi$ \\
$Q$ & per-type AMR carrying capacity &
  & \\
$\bm{a},\ a_j$ & AMR assignment vector; AMR assigned to job $j$ &
  \multicolumn{2}{@{}l}{\itshape Training and refinement} \\
$\sigma,\ \sigma_t$ & operation sequence; prefix committed up to step $t$ &
  $R,\ R(\sigma_{1:t})$ & dispatching rule; completion of prefix $\sigma_{1:t}$ by $R$ \\
$s^{\pi}_j,c^{\pi}_j$ & start / completion of pickup service of job $j$ &
  $C[\cdot]$ & shorthand for $C_{\max}(\Phi(\cdot))$ \\
$s^{\delta}_j,c^{\delta}_j$ & start / completion of delivery service of job $j$ &
  $A_t,\ \hat{A}_t$ & stepwise counterfactual advantage; standardized \\
$t_k(u,v)$ & realized travel time of AMR $k$ from $u$ to $v$ &
  $A^{\mathrm{lb}}_t$ & load-balance shaping advantage \\
$\Phi$ & collision-aware executor &
  $\lambda_{\mathrm{lb}},\ \lambda_{H}$ & shaping and entropy coefficients \\
$C_{\max},\ F_k$ & makespan; finish time of AMR $k$ incl.\ depot return &
  $\mathcal{H}_t$ & entropy of the masked policy at step $t$ \\
$W_d$ & number of physical waiting cells at dock $d$ &
  $\mathcal{L}(\theta)$ & training loss \\
\multicolumn{2}{@{}l}{\itshape Sequential decision process} &
  $B,\ E$ & batch size (episodes); number of training epochs \\
$s_t,\ a_t$ & state and action at decision epoch $t$ &
  $T_{\mathrm{val}},\ \kappa$ & validation interval; infeasibility penalty \\
$\omega$ & operation type, $\omega\in\{\textsc{pick},\textsc{deliver}\}$ &
  $I_{\mathrm{cf}},\ I_{\mathrm{ca}}$ & collision-free / collision-aware refinement budgets \\
$\Lambda(s_t)$ & feasible-action set (mask) in state $s_t$ &
  $\mathbf{1}[\cdot]$ & indicator function \\
$\hat{\Psi}$ & calibrated surrogate transition model &
  & \\
$\mathcal{E}_d$ & committed service events $(s_e,c_e,j_e,\omega_e)$ at dock $d$ &
  & \\
$\hat{b}_k,\ x_k$ & next-available time and position of AMR $k$ &
  & \\
$\rho$ & ready time of an AMR at a dock &
  & \\
$\hat{t}(u,v)$ & calibrated travel-time estimate \eqref{eq:travelcal} &
  & \\
$\hat{w}_d(\rho)$ & queue-aware dock-wait estimate \eqref{eq:waitcal} &
  & \\
$q_d(\rho)$ & services still unfinished at time $\rho$ at dock $d$ &
  & \\
$\alpha,\ \beta$ & travel-calibration scale and offset &
  & \\
$\varphi(\cdot)$ & dock-wait penalty table over queue depths &
  & \\
$\pi_{\theta}$ & scheduling policy with parameters $\theta$ &
  & \\
\bottomrule
\end{tabular}
\end{table*}

\subsection{System Description}

We consider the internal transfer problem of a cross-docking terminal
served by a fleet of autonomous mobile robots (AMRs). Table~\ref{tab:notation}
summarizes the notation used throughout the paper. The terminal floor
is modeled as a four-connected unit grid $\mathcal{W}$ in which a set of
\emph{inbound docks} $\mathcal{D}^{\mathrm{in}}$ is located along one
boundary and a set of \emph{outbound docks} (processing stations)
$\mathcal{D}^{\mathrm{out}}$ along the opposite boundary; the AMR depots
lie between them. Consistent with cross-docking practice, the floor is
free of static obstacles, so congestion arises from vehicle interference
and dock contention rather than from the layout itself. AMRs travel at
unit speed between adjacent cells, and the distance between two
positions $u,v$ is the Manhattan metric $\Delta(u,v)$.

A set of transfer jobs $\mathcal{J}=\{1,\dots,n\}$ must be moved across
the terminal by a fleet $\mathcal{M}=\{1,\dots,m\}$ of homogeneous AMRs.
Each job $j\in\mathcal{J}$ represents one unit load of a material type
$\tau_j\in\mathcal{T}$ that becomes available at its inbound dock
$o_j\in\mathcal{D}^{\mathrm{in}}$ at release time $r_j\ge 0$ and must be
delivered to a designated outbound dock
$d_j\in\mathcal{D}^{\mathrm{out}}$. Handling a unit load of type $\tau$
occupies a dock for a type-dependent service duration $p(\tau)$; we
write $p_j=p(\tau_j)$ and assume the loading service at the inbound dock
and the unloading/processing service at the outbound dock take the same
time, reflecting symmetric handling of a unit load. Each job therefore
consists of an ordered pair of operations: a \emph{pickup} operation
$\pi_j$ executed at $o_j$ and a \emph{delivery} operation $\delta_j$
executed at $d_j$, both performed by the same AMR (no transshipment).
An AMR may consolidate loads subject to a per-type carrying capacity
$Q$: at no time may it carry more than $Q$ units of any single material
type.

Docks are exclusive single-server resources: at most one AMR may be in
service at a dock at any time, and services are non-preemptive. An AMR
arriving at an occupied dock queues in a bounded set of waiting cells
adjacent to that dock; when the waiting line is saturated, the AMR must
hold at an upstream position and re-approach later. Finally, AMR motion
must be collision-free: at any unit time step, each grid cell may be
occupied by at most one AMR.

\subsection{Formal Statement}

A solution is a pair $(\bm{a},\sigma)$, where
$\bm{a}=(a_1,\dots,a_n)\in\mathcal{M}^{n}$ assigns each job to an AMR
and $\sigma$ is a total order over the $2n$ operations
$\{\pi_j,\delta_j\}_{j\in\mathcal{J}}$ that determines the sequence in
which operations are committed. Let $s^{\pi}_j,c^{\pi}_j$
(resp.\ $s^{\delta}_j,c^{\delta}_j$) denote the service start and
completion times of $\pi_j$ (resp.\ $\delta_j$) induced by executing
$(\bm{a},\sigma)$. A feasible execution satisfies:
\begin{align}
c^{\pi}_j &= s^{\pi}_j + p_j, \qquad
c^{\delta}_j = s^{\delta}_j + p_j,
&& \forall j\in\mathcal{J}, \label{eq:nonpreempt}\\
s^{\pi}_j &\ge r_j,
&& \forall j\in\mathcal{J}, \label{eq:release}\\
s^{\delta}_j &\ge c^{\pi}_j,
&& \forall j\in\mathcal{J}, \label{eq:precedence}
\end{align}
\begin{align}
\big|\{\,j : a_j{=}k,\ \tau_j{=}\tau,\ c^{\pi}_j \le u < c^{\delta}_j \}\big| &\le Q,
&& \forall k,\tau,u, \label{eq:capacity}\\
\big[s_e, c_e\big) \cap \big[s_{e'}, c_{e'}\big) &= \emptyset,
&& \forall e\ne e' \ \text{at the same dock}, \label{eq:exclusive}
\end{align}
where \eqref{eq:nonpreempt} enforces non-preemptive services,
\eqref{eq:release} release dates, \eqref{eq:precedence} pickup-before-%
delivery precedence, \eqref{eq:capacity} the per-type carrying capacity,
and \eqref{eq:exclusive} dock exclusivity over all service events $e$.
In addition, for any two operations $u\prec v$ executed consecutively by
the same AMR $k$,
\begin{equation}
s_v \;\ge\; c_u + t_k\big(\mathrm{loc}(u),\mathrm{loc}(v)\big),
\label{eq:travel}
\end{equation}
where $t_k(\cdot,\cdot)\ge\Delta(\cdot,\cdot)$ is the realized travel
time of AMR $k$, which may exceed the Manhattan distance due to
collision avoidance, queuing detours, and dock-clearing maneuvers.

The realized travel and waiting times are not fixed constants but are
determined by a high-fidelity \emph{executor} $\Phi$, which routes every
AMR leg with space--time $\mathrm{A}^{*}$ search over a shared
reservation table (at most one AMR per cell per time step), manages the
bounded dock waiting lines, and inserts upstream holding maneuvers when
lines are saturated. The problem is therefore a simulation-based
optimization: find
\begin{equation}
(\bm{a}^{*},\sigma^{*}) \;=\;
\arg\min_{(\bm{a},\sigma)} \; C_{\max}\big(\Phi(\bm{a},\sigma)\big),
\label{eq:objective}
\end{equation}
where the makespan $C_{\max}=\max_{k\in\mathcal{M}} F_k$ and $F_k$ is
the time at which AMR $k$ completes its final activity, including its
return trip to the depot. The problem generalizes both the capacitated
pickup-and-delivery problem and parallel-machine scheduling with
sequence-dependent setups, and is therefore NP-hard; the coupling to
collision-free multi-agent routing further compounds its difficulty
\cite{stern2019mapf}.

\subsection{Sequential Decision Formulation}
\label{sec:mdp}

We cast the construction of $(\bm{a},\sigma)$ as an episodic sequential
decision process with exactly $2n$ decision epochs, one per committed
operation. At epoch $t$, the state
\begin{equation}
s_t=\big(\{\bm{x}^{A}_{k}\},\, \{\bm{x}^{J}_{j}\},\,
\{\mathcal{E}_d\},\, \sigma_{t}\big)
\label{eq:state}
\end{equation}
collects the AMR states (position, next-available time, per-type
inventory), the job states (unpicked / onboard / delivered, carrier
identity), the set of committed dock service events
$\mathcal{E}_d=\{(s_e,c_e,j_e,\omega_e)\}$ at each dock $d$, and the
partial operation sequence $\sigma_t$.

An action is a triple
\begin{equation}
a_t=(\omega,k,j)\in\{\textsc{pick},\textsc{deliver}\}\times\mathcal{M}\times\mathcal{J},
\label{eq:action}
\end{equation}
giving a flat action space of size $2mn$ restricted by a feasibility
mask $\Lambda(s_t)$: $(\textsc{pick},k,j)$ is legal iff $j$ is unpicked
and AMR $k$ carries fewer than $Q$ units of type $\tau_j$;
$(\textsc{deliver},k,j)$ is legal iff $j$ is onboard AMR $k$. The mask
guarantees by construction that every complete episode yields a feasible
solution, so no repair mechanism is required.

Transitions follow a deterministic surrogate model $\hat{\Psi}$ that
advances the acting AMR's clock using a \emph{calibrated} travel-time
estimate $\hat{t}(\cdot,\cdot)$ and a queue-aware dock-wait estimate
$\hat{w}_d(\cdot)$ fitted offline against the executor $\Phi$
(Section~\ref{sec:calibration}); committing an action appends the
corresponding service window to $\mathcal{E}_d$. For a pickup, the
service start of job $j$ by AMR $k$ located at $x_k$ with next-available
time $\hat{b}_k$ is
\begin{equation}
s^{\pi}_j=\underbrace{\max\!\big(\hat{b}_k+\hat{t}(x_k,o_j),\, r_j\big)}_{\text{ready time }\rho}
\;+\;\hat{w}_{o_j}(\rho),
\label{eq:surrogate}
\end{equation}
and analogously for deliveries at $d_j$. After $2n$ steps the completed
solution is evaluated by the executor, and the learning objective is to
minimize the expected executed makespan
$\mathbb{E}_{\pi_\theta}\!\left[C_{\max}(\Phi(\bm{a},\sigma))\right]$;
the policy-gradient training procedure and its dispatch-rule baseline
are described in Section~\ref{sec:training}.

% =========================================================================

\section{Sparse-Graph Multi-Pointer Dispatcher}
\label{sec:architecture}

\subsection{Overview and Design Rationale}

Learning-to-dispatch architectures for shop scheduling encode
operations as nodes of a disjunctive graph and learn a construction
policy over it \cite{zhang2020l2d,song2023fjsp}. Transferring this
paradigm to AMR-served cross-docking raises a representation question:
how should the tight resource coupling---$m$ vehicles contending for
$|\mathcal{D}|$ single-server docks under collision-constrained
motion---enter the graph? A seemingly natural answer is to join every
pair of jobs that share a dock. We found this to be
counterproductive: with $n/|\mathcal{D}|$ jobs per dock, such
contention edges form dense cliques, and message passing over them
drives all job embeddings toward a common average (pairwise cosine
similarity above $0.99$ at initialization in our measurements),
leaving the policy unable to discriminate between candidate jobs
precisely at the opening decisions that dominate the makespan.
Over-smoothing under dense neighborhoods is a known failure mode of
deep graph networks \cite{li2018deeper,oono2020oversmoothing}; in this
problem it appears already at three message-passing rounds.

The proposed scheduler, hereafter \textsc{CD-GNN}, therefore keeps the
graph \emph{sparse and dynamic}: edges encode only scheduling
precedences that have actually been committed, while resource
congestion enters through calibrated, decision-grade node features.
The model comprises (i)~a job encoder over a dynamic
\emph{precedence graph} processed by residual graph isomorphism
network (GIN) layers \cite{xu2019gin}; (ii)~a lightweight fleet
encoder; and (iii)~a factored multi-pointer actor that scores each
masked (operation, AMR, job) triple \cite{kool2019attention}. The full
model is deliberately lean: $2.8\times10^{5}$ parameters
($2.5\times10^{5}$ actor, $3.3\times10^{4}$ critic) with hidden width
$h=128$ throughout, permutation-invariant in the job and AMR sets, and
applicable unchanged to any instance size.

% TODO: architecture overview figure
% \begin{figure*}[t]
%   \centering
%   \includegraphics[width=\textwidth]{figures/cdgnn_architecture}
%   \caption{Overview of \textsc{CD-GNN}. Job features (including
%   calibrated congestion estimates) are encoded and refined over the
%   dynamic precedence graph by residual GIN layers; AMR features are
%   encoded independently; the multi-pointer actor scores each feasible
%   (operation, AMR, job) triple under the feasibility mask.}
%   \label{fig:architecture}
% \end{figure*}

\subsection{Congestion-Aware State Representation}
\label{sec:features}

At every decision epoch the state \eqref{eq:state} is rendered into two
feature families, summarized in Table~\ref{tab:features}. All features
are normalized to comparable magnitudes by fixed scales.

\emph{Job nodes} ($\bm{x}^{J}_j\in\mathbb{R}^{16}$). Each job encodes
its service duration, lifecycle status (unpicked/onboard/delivered) and
carrier identity, its material type, and the coordinates of both its
supply position and its destination station, giving the network direct
geometric grounding. Resource congestion enters through four derived
features: the current availability delays of the job's inbound and
outbound docks, and the number of still-unpicked competitor jobs
committed to each of the two docks. A calibrated completion-time
estimate (Section~\ref{sec:calibration}) summarizes how quickly the job
could finish if served next, providing a value-like signal that the
network would otherwise need several layers to reconstruct.

\emph{AMR nodes} ($\bm{x}^{A}_k\in\mathbb{R}^{8}$). Each AMR encodes a
busy indicator and its time until available, per-type inventory
ratios, its coordinates, and the number of jobs currently assigned to
it, which exposes fleet-utilization imbalance to the policy.

\begin{table}[t]
\caption{Node Features of \textsc{CD-GNN}}
\label{tab:features}
\centering
\footnotesize
\begin{tabular}{@{}llc@{}}
\toprule
& Feature & Scale \\
\midrule
\multirow{12}{*}{\rotatebox{90}{Job $\bm{x}^{J}_j\!\in\!\mathbb{R}^{16}$}}
& existence flag & -- \\
& service duration $p_j$ & $1/25$ \\
& carrier index (0 if unpicked) & $1/m$ \\
& destination position $(x,y)$ ($\times2$) & $1/10$ \\
& supply position $(x,y)$ ($\times2$) & $1/10$ \\
& material type one-hot ($\times3$) & -- \\
& lifecycle status $\{0,\tfrac12,1\}$ & -- \\
& completion-time estimate & $1/500$ \\
& inbound dock availability delay & $1/100$ \\
& outbound dock availability delay & $1/100$ \\
& inbound dock committed queue & $1/m$ \\
& outbound dock committed queue & $1/m$ \\
\midrule
\multirow{7}{*}{\rotatebox{90}{AMR $\bm{x}^{A}_k\!\in\!\mathbb{R}^{8}$}}
& busy indicator & -- \\
& time until available & $1/50$ \\
& inventory ratio, per type ($\times3$) & $1/Q$ \\
& position $(x,y)$ ($\times2$) & $1/10$ \\
& assigned-job count & $1/10$ \\
\bottomrule
\end{tabular}
\end{table}

\subsection{Dynamic Precedence Graph and Job Encoder}
\label{sec:graph}

Job features pass through a linear map
$\bm{g}^{(0)}_j=\bm{W}_{\!J}\,\bm{x}^{J}_j\in\mathbb{R}^{h}$ and AMR
features through a linear encoder
$\bm{h}^{A}_k=\bm{W}_{\!A}\,\bm{x}^{A}_k\in\mathbb{R}^{h}$.

Jobs are coupled by the \emph{committed} structure of the partial
schedule. Following the adding-arc construction of disjunctive-graph
dispatching \cite{zhang2020l2d}, the graph
$\mathcal{G}_t=(\mathcal{J},E_t)$ starts empty and grows one or two
directed arcs per decision: when an operation of job $j$ is committed,
an arc is added from the previously scheduled job on the same AMR to
$j$, and from the previously scheduled job on the same outbound dock to
$j$. Each job thus aggregates information from its immediate
predecessors $\mathcal{P}(j)$ on both resources it occupies---the
vehicle sequence and the dock sequence---while uncommitted jobs
interact only through the shared congestion features of
Section~\ref{sec:features}, never through dense same-dock edges.

Message passing applies $L=3$ residual GIN layers with layer
normalization,
\begin{equation}
\bm{g}^{(l+1)}_j=\bm{g}^{(l)}_j+
\mathrm{LN}\!\left(\mathrm{MLP}^{(l)}\!\Big(
(1{+}\epsilon^{(l)})\,\bm{g}^{(l)}_j+
\sum_{i\in\mathcal{P}(j)}\!\bm{g}^{(l)}_i
\Big)\right)\!,
\label{eq:gin}
\end{equation}
retaining the injective sum aggregation of GIN \cite{xu2019gin}: with
at most two predecessors per node, message magnitudes stay bounded
irrespective of instance size, so the scale-invariance concerns that
arise with dense contention edges do not apply.

\subsection{Multi-Pointer Action Scoring and Value Estimation}
\label{sec:actor}

Let $\bm{e}_{\omega}\in\mathbb{R}^{h}$ be a learned embedding of the
operation type. The actor scores every candidate triple by
\begin{equation}
z_{\omega,k,j}=\mathrm{MLP}_{\pi}\!\big(\big[\,\bm{e}_{\omega}\,\|\,
\bm{g}^{(L)}_{j}\,\|\,\bm{h}^{A}_{k}\big]\big),
\label{eq:logit}
\end{equation}
and the policy is the masked softmax
\begin{equation}
\pi_{\theta}(a\,|\,s_t)=
\frac{\exp(z_{a})\,\mathbf{1}[a\in\Lambda(s_t)]}
{\sum_{a'\in\Lambda(s_t)}\exp(z_{a'})}.
\label{eq:policy}
\end{equation}
The critic estimates the state value from pooled entity summaries,
\begin{equation}
V_{\phi}(s_t)=\mathrm{MLP}_{V}\!\big(\big[\,
\overline{\bm{g}}^{\mathrm{(unf)}} \,\|\, \overline{\bm{h}}^{A}\big]\big),
\label{eq:critic}
\end{equation}
where $\overline{\bm{g}}^{\mathrm{(unf)}}$ averages the embeddings of
unfinished jobs only. A full episode requires exactly $2n$ forward
passes; each pass scores all $2mn$ actions in a single batched
evaluation of \eqref{eq:logit}.

% =========================================================================

\section{Surrogate Model Calibration}
\label{sec:calibration}

\subsection{The Two-Simulator Gap}

The policy constructs schedules on the surrogate $\hat{\Psi}$ of
Section~\ref{sec:mdp}, whereas solution quality is defined by the
executor $\Phi$ of \eqref{eq:objective}. An uncalibrated surrogate that
advances clocks by Manhattan distances and single-scalar dock
availabilities is systematically optimistic: it ignores collision
detours, dock-clearing maneuvers, waiting-line approach paths, and the
upstream holds triggered when waiting lines saturate. The congestion
features of Section~\ref{sec:features}---and therefore every decision
the policy makes---would then be conditioned on clocks that drift
further from reality as congestion grows, precisely in the regime where
scheduling decisions matter most. Rather than paying the computational
cost of embedding the full executor in the training loop, we close most
of this gap with two lightweight corrections fitted offline.

\subsection{Calibrated Travel and Dock-Wait Estimators}

\emph{Travel.} Realized leg times are modeled as an affine function of
the Manhattan distance, clamped from below by the metric itself:
\begin{equation}
\hat{t}(u,v)=\max\!\big(\Delta(u,v),\ \alpha\,\Delta(u,v)+\beta\big).
\label{eq:travelcal}
\end{equation}

\emph{Dock waiting.} The wait experienced by an AMR that becomes ready
at a dock $d$ at time $\rho$ is modeled as the committed availability
delay plus a penalty indexed by the queue depth at readiness,
\begin{equation}
\hat{w}_d(\rho)=\max\!\big(0,\ \mathrm{avail}_d-\rho\big)
+\varphi\!\big(\min(q_d(\rho),\,3)\big),
\label{eq:waitcal}
\end{equation}
where $q_d(\rho)=|\{e\in\mathcal{E}_d : c_e>\rho\}|$ counts committed
services unfinished at $\rho$, and
$\varphi:\{0,1,2,3^{+}\}\to\mathbb{R}_{\ge0}$ is a monotone lookup
table. The bounded table shape is deliberate: waiting-line detours and
upstream holds produce a cost cliff once the line saturates, a
nonlinearity that a linear congestion model cannot express.

\subsection{Fitting Procedure}

Calibration data are obtained by replaying complete schedules through
both simulators and aligning them per operation. Schedules are drawn
from heterogeneous sources---five dispatching-rule pairs and the
current learned policies---across instance sizes
$n\in\{10,20,40,60\}$, so that the fitted correction covers the state
distribution that policies actually visit. For every operation, the
raw surrogate supplies the Manhattan leg distance, the base wait, and
the queue depth at readiness; the executor's logged timeline is
partitioned, per AMR, into legs terminating at each service start,
within which movement entries sum to the realized travel and the
remainder is attributed to waiting. Executions containing infeasible
operations are discarded. The travel parameters $(\alpha,\beta)$ are
fitted by least squares and the penalties $\varphi(\cdot)$ as clipped
bin means with monotonicity enforced across depths.

In our environment, fitting on $10{,}360$ operation pairs from $159$
schedules yielded $\alpha=0.963$, $\beta=1.12$ ($R^{2}=0.91$) and
$\varphi=(0.04,\,0.04,\,1.66,\,3.57)$; the step between depths one and
two locates the waiting-line saturation cliff. The correction removes
the mean per-leg travel bias entirely ($+0.744\to0.000$ time units)
and halves the end-to-end makespan drift between surrogate and
executor (mean absolute error $24.3\to12.1$ on held-out schedules).
The calibration is a frozen artifact: it is fitted once before
training, adds no learned components to the training loop, and is
shared by \emph{all} policies and by the dispatching-rule baseline of
Section~\ref{sec:training}, so every method plans in the same
corrected world. After training converges, the fit can be re-estimated
from the new policies' schedules to confirm stability.

% =========================================================================

\section{Policy Optimization and Refinement}
\label{sec:training}

\subsection{Objective and Gradient Estimator}

The policy $\pi_\theta$ is trained to minimize the expected executed
makespan $\mathbb{E}_{\pi_\theta}[C_{\max}(\Phi(\bm{a},\sigma))]$ with
the REINFORCE estimator \cite{williams1992reinforce}. During training,
episodes are generated by sampling from the masked policy
\eqref{eq:policy} on the calibrated surrogate; at validation and
deployment, actions are selected greedily. The gradient quality of
REINFORCE hinges on its baseline; we use a counterfactual baseline
built from a strong dispatching rule and evaluated by the executor
itself.

\subsection{Counterfactual Dispatch-Rule Baseline}
\label{sec:baseline}

Let $R$ denote a deterministic dispatching heuristic (job rule and AMR
rule; we use earliest-completion for both) and let
$R(\sigma_{1:t})$ denote the complete solution obtained by executing
the partial sequence $\sigma_{1:t}$ and completing all remaining
operations with $R$. Define the shorthand
$C[\cdot]=C_{\max}(\Phi(\cdot))$. The advantage of the action taken at
step $t$ is the improvement of committing that action and then
following the rule, relative to following the rule immediately:
\begin{equation}
A_t \;=\; C\big[R(\sigma_{1:t-1})\big]\;-\;
C\big[R(\sigma_{1:t-1}\oplus a_t)\big].
\label{eq:advantage}
\end{equation}
Both completions are evaluated by the collision-aware executor, so the
learning signal is grounded in true solution quality rather than in
surrogate clocks. Unlike episode-level rollout baselines
\cite{rennie2017self,kool2019attention}, \eqref{eq:advantage} assigns
credit per decision: an early pickup that congests a dock is penalized
at the step where it is committed, not diluted across the episode.
The price is $\mathcal{O}(n)$ executor evaluations per episode, which
dominates training cost; we consider the sharper credit assignment a
favorable trade at the instance sizes studied here, and
Section~\ref{sec:training-details} lists a cheaper actor--critic
alternative.

\subsection{Loss Function}

Advantages are standardized to zero mean and unit variance across all
steps of a batch of $B$ episodes. The per-batch loss combines the
policy gradient with a load-balancing shaping term and an entropy
bonus:
\begin{equation}
\mathcal{L}(\theta)=
-\frac{1}{B}\sum_{b=1}^{B}\sum_{t=1}^{2n}
\Big[\log\pi_\theta(a_t^b|s_t^b)\,
\big(\hat{A}_t^{\,b}+\lambda_{\mathrm{lb}}A_t^{\mathrm{lb},b}\big)
+\lambda_{H}\,\mathcal{H}_t^{b}\Big],
\label{eq:loss}
\end{equation}
where $\hat{A}_t$ is the standardized advantage,
$\mathcal{H}_t$ is the entropy of the masked policy over feasible
actions, and the shaping term
$A^{\mathrm{lb}}_t=\min_{k'}c_{k'}-c_{k}$ (for pickups; zero for
deliveries) compares the assignment count $c_k$ of the chosen AMR with
the least-loaded AMR, discouraging degenerate fleet utilization early
in training \cite{ng1999shaping}. We use
$\lambda_{\mathrm{lb}}=0.1$ and $\lambda_{H}=0.01$.

\begin{algorithm}[t]
\caption{Training with counterfactual dispatch baseline}
\label{alg:training}
\begin{algorithmic}[1]
\Require policy $\pi_\theta$, rule $R$, executor $\Phi$, surrogate $\hat{\Psi}$
\State fit calibration $(\alpha,\beta,\varphi)$ offline; freeze (Sec.~\ref{sec:calibration})
\For{$\mathrm{epoch}=1,\dots,E$}
  \For{$b=1,\dots,B$}
    \State sample instance; roll out $\sigma^b\!\sim\!\pi_\theta$ on $\hat{\Psi}$
    \For{$t=1,\dots,2n$}
      \State $A^b_t \gets C[R(\sigma^b_{1:t-1})]-C[R(\sigma^b_{1:t-1}\oplus a^b_t)]$
    \EndFor
  \EndFor
  \State standardize $\{A^b_t\}$ over the batch
  \State update $\theta$ by Adam on $\mathcal{L}(\theta)$ in \eqref{eq:loss}; clip $\lVert\nabla\rVert\le1$
  \If{$\mathrm{epoch}\bmod T_{\mathrm{val}}=0$}
    \State greedy-decode held-out instances; score by
    $C_{\max}+\kappa\cdot(\text{infeasible ops})$; keep best
  \EndIf
\EndFor
\end{algorithmic}
\end{algorithm}

\subsection{Training Protocol}
\label{sec:training-details}

Algorithm~\ref{alg:training} summarizes the procedure. Only the actor
parameters (encoders, GIN layers, operation embedding, and actor head)
receive gradients under REINFORCE; the critic \eqref{eq:critic} is
reserved for an optional proximal-policy-optimization variant
\cite{schulman2017ppo} that replaces \eqref{eq:advantage} with
critic-based advantages when the executor evaluations of the stepwise
baseline are prohibitive. We train with Adam \cite{kingma2015adam}
(learning rate $3\times10^{-4}$, batch $B=16$, gradient-norm clip
$1.0$) for $2{,}000$ epochs. Every $T_{\mathrm{val}}=50$ epochs the
policy is evaluated by greedy decoding on a fixed held-out validation
set, disjoint from the training pool, and scored by the executed
makespan plus a penalty of $\kappa=10^{3}$ per infeasible operation;
the checkpoint with the best validation score is retained. Computing
the entropy in \eqref{eq:loss} over the feasible-action support only
keeps the bonus comparable as the mask tightens toward the end of an
episode.

\subsection{Post-Processing Local Search}
\label{sec:localsearch}

At deployment, the greedy-decoded schedule is refined by a two-stage
first-improvement local search that reuses the neighborhood operators
of our genetic-algorithm baseline, so that quality differences between
methods are attributable to the constructed solutions rather than to
the refinement. A move either swaps two operations in $\sigma$, at
least one of which belongs to the AMR that currently determines the
makespan, or relocates a contiguous block of same-AMR, same-type
operations; the AMR assignment $\bm{a}$ is left unchanged. The first
stage performs $I_{\mathrm{cf}}=10^{3}$ trials under the fast
collision-free evaluator to explore broadly; the second performs
$I_{\mathrm{ca}}=10^{2}$ trials under the full executor $\Phi$ to
polish under true routing. The refinement adds approximately two
seconds per instance and is entirely optional: all learned-policy
results are reported both with and without it.

% =========================================================================
% Suggested BibTeX entries (replace keys above with your bibliography):
%
% @inproceedings{xu2019gin, Xu, Hu, Leskovec, Jegelka, "How Powerful Are
%   Graph Neural Networks?", ICLR 2019}
% @inproceedings{zhang2020l2d, Zhang et al., "Learning to Dispatch for
%   Job Shop Scheduling via Deep Reinforcement Learning", NeurIPS 2020}
% @article{song2023fjsp, Song et al., "Flexible Job-Shop Scheduling via
%   Graph Neural Network and Deep Reinforcement Learning", IEEE TII 2023}
% @inproceedings{li2018deeper, Li, Han, Wu, "Deeper Insights into Graph
%   Convolutional Networks for Semi-Supervised Learning", AAAI 2018}
% @inproceedings{oono2020oversmoothing, Oono, Suzuki, "Graph Neural
%   Networks Exponentially Lose Expressive Power for Node
%   Classification", ICLR 2020}
% @inproceedings{stern2019mapf, Stern et al., "Multi-Agent Pathfinding:
%   Definitions, Variants, and Benchmarks", SoCS 2019}
% @article{williams1992reinforce, Williams, "Simple Statistical
%   Gradient-Following Algorithms for Connectionist Reinforcement
%   Learning", Machine Learning 8, 1992}
% @inproceedings{rennie2017self, Rennie et al., "Self-Critical Sequence
%   Training for Image Captioning", CVPR 2017}
% @inproceedings{kool2019attention, Kool, van Hoof, Welling, "Attention,
%   Learn to Solve Routing Problems!", ICLR 2019}
% @inproceedings{ng1999shaping, Ng, Harada, Russell, "Policy Invariance
%   Under Reward Transformations: Theory and Application to Reward
%   Shaping", ICML 1999}
% @article{schulman2017ppo, Schulman et al., "Proximal Policy
%   Optimization Algorithms", arXiv:1707.06347}
% @inproceedings{kingma2015adam, Kingma, Ba, "Adam: A Method for
%   Stochastic Optimization", ICLR 2015}
% =========================================================================
 inside an IEEEtran doc.
% Requires in the preamble:
%   \usepackage{amsmath,amssymb,bm}
%   \usepackage{booktabs,array,multirow,graphicx}
%   \usepackage{algorithm}
%   \usepackage{algpseudocode}
% Citation keys are placeholders -- suggested BibTeX entries are listed in
% comments at the end. "CD-GNN" (Cross-Docking GNN) is a placeholder model
% name; rename globally if you prefer.
% Sections: Problem Formulation, Architecture, Surrogate Calibration,
% Policy Optimization (incl. post-processing local search).
% =========================================================================

\section{Problem Formulation}
\label{sec:problem}

\begin{table*}[t]
\caption{Summary of Notation}
\label{tab:notation}
\centering
\footnotesize
\begin{tabular}{@{}l p{0.295\textwidth} @{\qquad} l p{0.295\textwidth}@{}}
\toprule
Symbol & Meaning & Symbol & Meaning \\
\midrule
\multicolumn{2}{@{}l}{\itshape Problem setting} &
\multicolumn{2}{@{}l}{\itshape Architecture} \\
$\mathcal{W}$ & grid workspace (four-connected unit cells) &
  $\bm{x}^{A}_k,\bm{x}^{J}_j$ & AMR / job input features ($\mathbb{R}^{8},\mathbb{R}^{16}$) \\
$\Delta(u,v)$ & Manhattan distance between positions $u,v$ &
  $h$ & hidden width of all embeddings ($h{=}128$) \\
$\mathcal{J},\ n$ & job set; number of jobs &
  $\bm{h}^{A}_k$ & AMR embedding \\
$\mathcal{M},\ m$ & AMR fleet; fleet size &
  $\bm{g}^{(l)}_j$ & job embedding after $l$ GIN layers \\
$\mathcal{T}$ & set of material types &
  $\mathcal{G}_t=(\mathcal{J},E_t)$ & dynamic precedence graph at step $t$ \\
$\mathcal{D}^{\mathrm{in}},\mathcal{D}^{\mathrm{out}},\mathcal{D}$ & inbound docks; outbound docks; their union &
  $\mathcal{P}(j)$ & committed predecessors of job $j$ in $\mathcal{G}_t$ \\
$\tau_j$ & material type of job $j$ &
  $\epsilon^{(l)},\ L$ & learnable GIN self-weight; number of GIN layers \\
$o_j,\ d_j$ & inbound and outbound dock of job $j$ &
  $\mathrm{LN}$ & layer normalization \\
$r_j$ & release time of job $j$ &
  $\bm{e}_{\omega}$ & learned operation-type embedding \\
$p_j=p(\tau_j)$ & dock service duration of job $j$ &
  $z_{\omega,k,j}$ & logit of action $(\omega,k,j)$ \\
$\pi_j,\ \delta_j$ & pickup and delivery operations of job $j$ &
  $V_{\phi}$ & critic (state-value) network, parameters $\phi$ \\
$Q$ & per-type AMR carrying capacity &
  & \\
$\bm{a},\ a_j$ & AMR assignment vector; AMR assigned to job $j$ &
  \multicolumn{2}{@{}l}{\itshape Training and refinement} \\
$\sigma,\ \sigma_t$ & operation sequence; prefix committed up to step $t$ &
  $R,\ R(\sigma_{1:t})$ & dispatching rule; completion of prefix $\sigma_{1:t}$ by $R$ \\
$s^{\pi}_j,c^{\pi}_j$ & start / completion of pickup service of job $j$ &
  $C[\cdot]$ & shorthand for $C_{\max}(\Phi(\cdot))$ \\
$s^{\delta}_j,c^{\delta}_j$ & start / completion of delivery service of job $j$ &
  $A_t,\ \hat{A}_t$ & stepwise counterfactual advantage; standardized \\
$t_k(u,v)$ & realized travel time of AMR $k$ from $u$ to $v$ &
  $A^{\mathrm{lb}}_t$ & load-balance shaping advantage \\
$\Phi$ & collision-aware executor &
  $\lambda_{\mathrm{lb}},\ \lambda_{H}$ & shaping and entropy coefficients \\
$C_{\max},\ F_k$ & makespan; finish time of AMR $k$ incl.\ depot return &
  $\mathcal{H}_t$ & entropy of the masked policy at step $t$ \\
$W_d$ & number of physical waiting cells at dock $d$ &
  $\mathcal{L}(\theta)$ & training loss \\
\multicolumn{2}{@{}l}{\itshape Sequential decision process} &
  $B,\ E$ & batch size (episodes); number of training epochs \\
$s_t,\ a_t$ & state and action at decision epoch $t$ &
  $T_{\mathrm{val}},\ \kappa$ & validation interval; infeasibility penalty \\
$\omega$ & operation type, $\omega\in\{\textsc{pick},\textsc{deliver}\}$ &
  $I_{\mathrm{cf}},\ I_{\mathrm{ca}}$ & collision-free / collision-aware refinement budgets \\
$\Lambda(s_t)$ & feasible-action set (mask) in state $s_t$ &
  $\mathbf{1}[\cdot]$ & indicator function \\
$\hat{\Psi}$ & calibrated surrogate transition model &
  & \\
$\mathcal{E}_d$ & committed service events $(s_e,c_e,j_e,\omega_e)$ at dock $d$ &
  & \\
$\hat{b}_k,\ x_k$ & next-available time and position of AMR $k$ &
  & \\
$\rho$ & ready time of an AMR at a dock &
  & \\
$\hat{t}(u,v)$ & calibrated travel-time estimate \eqref{eq:travelcal} &
  & \\
$\hat{w}_d(\rho)$ & queue-aware dock-wait estimate \eqref{eq:waitcal} &
  & \\
$q_d(\rho)$ & services still unfinished at time $\rho$ at dock $d$ &
  & \\
$\alpha,\ \beta$ & travel-calibration scale and offset &
  & \\
$\varphi(\cdot)$ & dock-wait penalty table over queue depths &
  & \\
$\pi_{\theta}$ & scheduling policy with parameters $\theta$ &
  & \\
\bottomrule
\end{tabular}
\end{table*}

\subsection{System Description}

We consider the internal transfer problem of a cross-docking terminal
served by a fleet of autonomous mobile robots (AMRs). Table~\ref{tab:notation}
summarizes the notation used throughout the paper. The terminal floor
is modeled as a four-connected unit grid $\mathcal{W}$ in which a set of
\emph{inbound docks} $\mathcal{D}^{\mathrm{in}}$ is located along one
boundary and a set of \emph{outbound docks} (processing stations)
$\mathcal{D}^{\mathrm{out}}$ along the opposite boundary; the AMR depots
lie between them. Consistent with cross-docking practice, the floor is
free of static obstacles, so congestion arises from vehicle interference
and dock contention rather than from the layout itself. AMRs travel at
unit speed between adjacent cells, and the distance between two
positions $u,v$ is the Manhattan metric $\Delta(u,v)$.

A set of transfer jobs $\mathcal{J}=\{1,\dots,n\}$ must be moved across
the terminal by a fleet $\mathcal{M}=\{1,\dots,m\}$ of homogeneous AMRs.
Each job $j\in\mathcal{J}$ represents one unit load of a material type
$\tau_j\in\mathcal{T}$ that becomes available at its inbound dock
$o_j\in\mathcal{D}^{\mathrm{in}}$ at release time $r_j\ge 0$ and must be
delivered to a designated outbound dock
$d_j\in\mathcal{D}^{\mathrm{out}}$. Handling a unit load of type $\tau$
occupies a dock for a type-dependent service duration $p(\tau)$; we
write $p_j=p(\tau_j)$ and assume the loading service at the inbound dock
and the unloading/processing service at the outbound dock take the same
time, reflecting symmetric handling of a unit load. Each job therefore
consists of an ordered pair of operations: a \emph{pickup} operation
$\pi_j$ executed at $o_j$ and a \emph{delivery} operation $\delta_j$
executed at $d_j$, both performed by the same AMR (no transshipment).
An AMR may consolidate loads subject to a per-type carrying capacity
$Q$: at no time may it carry more than $Q$ units of any single material
type.

Docks are exclusive single-server resources: at most one AMR may be in
service at a dock at any time, and services are non-preemptive. An AMR
arriving at an occupied dock queues in a bounded set of waiting cells
adjacent to that dock; when the waiting line is saturated, the AMR must
hold at an upstream position and re-approach later. Finally, AMR motion
must be collision-free: at any unit time step, each grid cell may be
occupied by at most one AMR.

\subsection{Formal Statement}

A solution is a pair $(\bm{a},\sigma)$, where
$\bm{a}=(a_1,\dots,a_n)\in\mathcal{M}^{n}$ assigns each job to an AMR
and $\sigma$ is a total order over the $2n$ operations
$\{\pi_j,\delta_j\}_{j\in\mathcal{J}}$ that determines the sequence in
which operations are committed. Let $s^{\pi}_j,c^{\pi}_j$
(resp.\ $s^{\delta}_j,c^{\delta}_j$) denote the service start and
completion times of $\pi_j$ (resp.\ $\delta_j$) induced by executing
$(\bm{a},\sigma)$. A feasible execution satisfies:
\begin{align}
c^{\pi}_j &= s^{\pi}_j + p_j, \qquad
c^{\delta}_j = s^{\delta}_j + p_j,
&& \forall j\in\mathcal{J}, \label{eq:nonpreempt}\\
s^{\pi}_j &\ge r_j,
&& \forall j\in\mathcal{J}, \label{eq:release}\\
s^{\delta}_j &\ge c^{\pi}_j,
&& \forall j\in\mathcal{J}, \label{eq:precedence}
\end{align}
\begin{align}
\big|\{\,j : a_j{=}k,\ \tau_j{=}\tau,\ c^{\pi}_j \le u < c^{\delta}_j \}\big| &\le Q,
&& \forall k,\tau,u, \label{eq:capacity}\\
\big[s_e, c_e\big) \cap \big[s_{e'}, c_{e'}\big) &= \emptyset,
&& \forall e\ne e' \ \text{at the same dock}, \label{eq:exclusive}
\end{align}
where \eqref{eq:nonpreempt} enforces non-preemptive services,
\eqref{eq:release} release dates, \eqref{eq:precedence} pickup-before-%
delivery precedence, \eqref{eq:capacity} the per-type carrying capacity,
and \eqref{eq:exclusive} dock exclusivity over all service events $e$.
In addition, for any two operations $u\prec v$ executed consecutively by
the same AMR $k$,
\begin{equation}
s_v \;\ge\; c_u + t_k\big(\mathrm{loc}(u),\mathrm{loc}(v)\big),
\label{eq:travel}
\end{equation}
where $t_k(\cdot,\cdot)\ge\Delta(\cdot,\cdot)$ is the realized travel
time of AMR $k$, which may exceed the Manhattan distance due to
collision avoidance, queuing detours, and dock-clearing maneuvers.

The realized travel and waiting times are not fixed constants but are
determined by a high-fidelity \emph{executor} $\Phi$, which routes every
AMR leg with space--time $\mathrm{A}^{*}$ search over a shared
reservation table (at most one AMR per cell per time step), manages the
bounded dock waiting lines, and inserts upstream holding maneuvers when
lines are saturated. The problem is therefore a simulation-based
optimization: find
\begin{equation}
(\bm{a}^{*},\sigma^{*}) \;=\;
\arg\min_{(\bm{a},\sigma)} \; C_{\max}\big(\Phi(\bm{a},\sigma)\big),
\label{eq:objective}
\end{equation}
where the makespan $C_{\max}=\max_{k\in\mathcal{M}} F_k$ and $F_k$ is
the time at which AMR $k$ completes its final activity, including its
return trip to the depot. The problem generalizes both the capacitated
pickup-and-delivery problem and parallel-machine scheduling with
sequence-dependent setups, and is therefore NP-hard; the coupling to
collision-free multi-agent routing further compounds its difficulty
\cite{stern2019mapf}.

\subsection{Sequential Decision Formulation}
\label{sec:mdp}

We cast the construction of $(\bm{a},\sigma)$ as an episodic sequential
decision process with exactly $2n$ decision epochs, one per committed
operation. At epoch $t$, the state
\begin{equation}
s_t=\big(\{\bm{x}^{A}_{k}\},\, \{\bm{x}^{J}_{j}\},\,
\{\mathcal{E}_d\},\, \sigma_{t}\big)
\label{eq:state}
\end{equation}
collects the AMR states (position, next-available time, per-type
inventory), the job states (unpicked / onboard / delivered, carrier
identity), the set of committed dock service events
$\mathcal{E}_d=\{(s_e,c_e,j_e,\omega_e)\}$ at each dock $d$, and the
partial operation sequence $\sigma_t$.

An action is a triple
\begin{equation}
a_t=(\omega,k,j)\in\{\textsc{pick},\textsc{deliver}\}\times\mathcal{M}\times\mathcal{J},
\label{eq:action}
\end{equation}
giving a flat action space of size $2mn$ restricted by a feasibility
mask $\Lambda(s_t)$: $(\textsc{pick},k,j)$ is legal iff $j$ is unpicked
and AMR $k$ carries fewer than $Q$ units of type $\tau_j$;
$(\textsc{deliver},k,j)$ is legal iff $j$ is onboard AMR $k$. The mask
guarantees by construction that every complete episode yields a feasible
solution, so no repair mechanism is required.

Transitions follow a deterministic surrogate model $\hat{\Psi}$ that
advances the acting AMR's clock using a \emph{calibrated} travel-time
estimate $\hat{t}(\cdot,\cdot)$ and a queue-aware dock-wait estimate
$\hat{w}_d(\cdot)$ fitted offline against the executor $\Phi$
(Section~\ref{sec:calibration}); committing an action appends the
corresponding service window to $\mathcal{E}_d$. For a pickup, the
service start of job $j$ by AMR $k$ located at $x_k$ with next-available
time $\hat{b}_k$ is
\begin{equation}
s^{\pi}_j=\underbrace{\max\!\big(\hat{b}_k+\hat{t}(x_k,o_j),\, r_j\big)}_{\text{ready time }\rho}
\;+\;\hat{w}_{o_j}(\rho),
\label{eq:surrogate}
\end{equation}
and analogously for deliveries at $d_j$. After $2n$ steps the completed
solution is evaluated by the executor, and the learning objective is to
minimize the expected executed makespan
$\mathbb{E}_{\pi_\theta}\!\left[C_{\max}(\Phi(\bm{a},\sigma))\right]$;
the policy-gradient training procedure and its dispatch-rule baseline
are described in Section~\ref{sec:training}.

% =========================================================================

\section{Sparse-Graph Multi-Pointer Dispatcher}
\label{sec:architecture}

\subsection{Overview and Design Rationale}

Learning-to-dispatch architectures for shop scheduling encode
operations as nodes of a disjunctive graph and learn a construction
policy over it \cite{zhang2020l2d,song2023fjsp}. Transferring this
paradigm to AMR-served cross-docking raises a representation question:
how should the tight resource coupling---$m$ vehicles contending for
$|\mathcal{D}|$ single-server docks under collision-constrained
motion---enter the graph? A seemingly natural answer is to join every
pair of jobs that share a dock. We found this to be
counterproductive: with $n/|\mathcal{D}|$ jobs per dock, such
contention edges form dense cliques, and message passing over them
drives all job embeddings toward a common average (pairwise cosine
similarity above $0.99$ at initialization in our measurements),
leaving the policy unable to discriminate between candidate jobs
precisely at the opening decisions that dominate the makespan.
Over-smoothing under dense neighborhoods is a known failure mode of
deep graph networks \cite{li2018deeper,oono2020oversmoothing}; in this
problem it appears already at three message-passing rounds.

The proposed scheduler, hereafter \textsc{CD-GNN}, therefore keeps the
graph \emph{sparse and dynamic}: edges encode only scheduling
precedences that have actually been committed, while resource
congestion enters through calibrated, decision-grade node features.
The model comprises (i)~a job encoder over a dynamic
\emph{precedence graph} processed by residual graph isomorphism
network (GIN) layers \cite{xu2019gin}; (ii)~a lightweight fleet
encoder; and (iii)~a factored multi-pointer actor that scores each
masked (operation, AMR, job) triple \cite{kool2019attention}. The full
model is deliberately lean: $2.8\times10^{5}$ parameters
($2.5\times10^{5}$ actor, $3.3\times10^{4}$ critic) with hidden width
$h=128$ throughout, permutation-invariant in the job and AMR sets, and
applicable unchanged to any instance size.

% TODO: architecture overview figure
% \begin{figure*}[t]
%   \centering
%   \includegraphics[width=\textwidth]{figures/cdgnn_architecture}
%   \caption{Overview of \textsc{CD-GNN}. Job features (including
%   calibrated congestion estimates) are encoded and refined over the
%   dynamic precedence graph by residual GIN layers; AMR features are
%   encoded independently; the multi-pointer actor scores each feasible
%   (operation, AMR, job) triple under the feasibility mask.}
%   \label{fig:architecture}
% \end{figure*}

\subsection{Congestion-Aware State Representation}
\label{sec:features}

At every decision epoch the state \eqref{eq:state} is rendered into two
feature families, summarized in Table~\ref{tab:features}. All features
are normalized to comparable magnitudes by fixed scales.

\emph{Job nodes} ($\bm{x}^{J}_j\in\mathbb{R}^{16}$). Each job encodes
its service duration, lifecycle status (unpicked/onboard/delivered) and
carrier identity, its material type, and the coordinates of both its
supply position and its destination station, giving the network direct
geometric grounding. Resource congestion enters through four derived
features: the current availability delays of the job's inbound and
outbound docks, and the number of still-unpicked competitor jobs
committed to each of the two docks. A calibrated completion-time
estimate (Section~\ref{sec:calibration}) summarizes how quickly the job
could finish if served next, providing a value-like signal that the
network would otherwise need several layers to reconstruct.

\emph{AMR nodes} ($\bm{x}^{A}_k\in\mathbb{R}^{8}$). Each AMR encodes a
busy indicator and its time until available, per-type inventory
ratios, its coordinates, and the number of jobs currently assigned to
it, which exposes fleet-utilization imbalance to the policy.

\begin{table}[t]
\caption{Node Features of \textsc{CD-GNN}}
\label{tab:features}
\centering
\footnotesize
\begin{tabular}{@{}llc@{}}
\toprule
& Feature & Scale \\
\midrule
\multirow{12}{*}{\rotatebox{90}{Job $\bm{x}^{J}_j\!\in\!\mathbb{R}^{16}$}}
& existence flag & -- \\
& service duration $p_j$ & $1/25$ \\
& carrier index (0 if unpicked) & $1/m$ \\
& destination position $(x,y)$ ($\times2$) & $1/10$ \\
& supply position $(x,y)$ ($\times2$) & $1/10$ \\
& material type one-hot ($\times3$) & -- \\
& lifecycle status $\{0,\tfrac12,1\}$ & -- \\
& completion-time estimate & $1/500$ \\
& inbound dock availability delay & $1/100$ \\
& outbound dock availability delay & $1/100$ \\
& inbound dock committed queue & $1/m$ \\
& outbound dock committed queue & $1/m$ \\
\midrule
\multirow{7}{*}{\rotatebox{90}{AMR $\bm{x}^{A}_k\!\in\!\mathbb{R}^{8}$}}
& busy indicator & -- \\
& time until available & $1/50$ \\
& inventory ratio, per type ($\times3$) & $1/Q$ \\
& position $(x,y)$ ($\times2$) & $1/10$ \\
& assigned-job count & $1/10$ \\
\bottomrule
\end{tabular}
\end{table}

\subsection{Dynamic Precedence Graph and Job Encoder}
\label{sec:graph}

Job features pass through a linear map
$\bm{g}^{(0)}_j=\bm{W}_{\!J}\,\bm{x}^{J}_j\in\mathbb{R}^{h}$ and AMR
features through a linear encoder
$\bm{h}^{A}_k=\bm{W}_{\!A}\,\bm{x}^{A}_k\in\mathbb{R}^{h}$.

Jobs are coupled by the \emph{committed} structure of the partial
schedule. Following the adding-arc construction of disjunctive-graph
dispatching \cite{zhang2020l2d}, the graph
$\mathcal{G}_t=(\mathcal{J},E_t)$ starts empty and grows one or two
directed arcs per decision: when an operation of job $j$ is committed,
an arc is added from the previously scheduled job on the same AMR to
$j$, and from the previously scheduled job on the same outbound dock to
$j$. Each job thus aggregates information from its immediate
predecessors $\mathcal{P}(j)$ on both resources it occupies---the
vehicle sequence and the dock sequence---while uncommitted jobs
interact only through the shared congestion features of
Section~\ref{sec:features}, never through dense same-dock edges.

Message passing applies $L=3$ residual GIN layers with layer
normalization,
\begin{equation}
\bm{g}^{(l+1)}_j=\bm{g}^{(l)}_j+
\mathrm{LN}\!\left(\mathrm{MLP}^{(l)}\!\Big(
(1{+}\epsilon^{(l)})\,\bm{g}^{(l)}_j+
\sum_{i\in\mathcal{P}(j)}\!\bm{g}^{(l)}_i
\Big)\right)\!,
\label{eq:gin}
\end{equation}
retaining the injective sum aggregation of GIN \cite{xu2019gin}: with
at most two predecessors per node, message magnitudes stay bounded
irrespective of instance size, so the scale-invariance concerns that
arise with dense contention edges do not apply.

\subsection{Multi-Pointer Action Scoring and Value Estimation}
\label{sec:actor}

Let $\bm{e}_{\omega}\in\mathbb{R}^{h}$ be a learned embedding of the
operation type. The actor scores every candidate triple by
\begin{equation}
z_{\omega,k,j}=\mathrm{MLP}_{\pi}\!\big(\big[\,\bm{e}_{\omega}\,\|\,
\bm{g}^{(L)}_{j}\,\|\,\bm{h}^{A}_{k}\big]\big),
\label{eq:logit}
\end{equation}
and the policy is the masked softmax
\begin{equation}
\pi_{\theta}(a\,|\,s_t)=
\frac{\exp(z_{a})\,\mathbf{1}[a\in\Lambda(s_t)]}
{\sum_{a'\in\Lambda(s_t)}\exp(z_{a'})}.
\label{eq:policy}
\end{equation}
The critic estimates the state value from pooled entity summaries,
\begin{equation}
V_{\phi}(s_t)=\mathrm{MLP}_{V}\!\big(\big[\,
\overline{\bm{g}}^{\mathrm{(unf)}} \,\|\, \overline{\bm{h}}^{A}\big]\big),
\label{eq:critic}
\end{equation}
where $\overline{\bm{g}}^{\mathrm{(unf)}}$ averages the embeddings of
unfinished jobs only. A full episode requires exactly $2n$ forward
passes; each pass scores all $2mn$ actions in a single batched
evaluation of \eqref{eq:logit}.

% =========================================================================

\section{Surrogate Model Calibration}
\label{sec:calibration}

\subsection{The Two-Simulator Gap}

The policy constructs schedules on the surrogate $\hat{\Psi}$ of
Section~\ref{sec:mdp}, whereas solution quality is defined by the
executor $\Phi$ of \eqref{eq:objective}. An uncalibrated surrogate that
advances clocks by Manhattan distances and single-scalar dock
availabilities is systematically optimistic: it ignores collision
detours, dock-clearing maneuvers, waiting-line approach paths, and the
upstream holds triggered when waiting lines saturate. The congestion
features of Section~\ref{sec:features}---and therefore every decision
the policy makes---would then be conditioned on clocks that drift
further from reality as congestion grows, precisely in the regime where
scheduling decisions matter most. Rather than paying the computational
cost of embedding the full executor in the training loop, we close most
of this gap with two lightweight corrections fitted offline.

\subsection{Calibrated Travel and Dock-Wait Estimators}

\emph{Travel.} Realized leg times are modeled as an affine function of
the Manhattan distance, clamped from below by the metric itself:
\begin{equation}
\hat{t}(u,v)=\max\!\big(\Delta(u,v),\ \alpha\,\Delta(u,v)+\beta\big).
\label{eq:travelcal}
\end{equation}

\emph{Dock waiting.} The wait experienced by an AMR that becomes ready
at a dock $d$ at time $\rho$ is modeled as the committed availability
delay plus a penalty indexed by the queue depth at readiness,
\begin{equation}
\hat{w}_d(\rho)=\max\!\big(0,\ \mathrm{avail}_d-\rho\big)
+\varphi\!\big(\min(q_d(\rho),\,3)\big),
\label{eq:waitcal}
\end{equation}
where $q_d(\rho)=|\{e\in\mathcal{E}_d : c_e>\rho\}|$ counts committed
services unfinished at $\rho$, and
$\varphi:\{0,1,2,3^{+}\}\to\mathbb{R}_{\ge0}$ is a monotone lookup
table. The bounded table shape is deliberate: waiting-line detours and
upstream holds produce a cost cliff once the line saturates, a
nonlinearity that a linear congestion model cannot express.

\subsection{Fitting Procedure}

Calibration data are obtained by replaying complete schedules through
both simulators and aligning them per operation. Schedules are drawn
from heterogeneous sources---five dispatching-rule pairs and the
current learned policies---across instance sizes
$n\in\{10,20,40,60\}$, so that the fitted correction covers the state
distribution that policies actually visit. For every operation, the
raw surrogate supplies the Manhattan leg distance, the base wait, and
the queue depth at readiness; the executor's logged timeline is
partitioned, per AMR, into legs terminating at each service start,
within which movement entries sum to the realized travel and the
remainder is attributed to waiting. Executions containing infeasible
operations are discarded. The travel parameters $(\alpha,\beta)$ are
fitted by least squares and the penalties $\varphi(\cdot)$ as clipped
bin means with monotonicity enforced across depths.

In our environment, fitting on $10{,}360$ operation pairs from $159$
schedules yielded $\alpha=0.963$, $\beta=1.12$ ($R^{2}=0.91$) and
$\varphi=(0.04,\,0.04,\,1.66,\,3.57)$; the step between depths one and
two locates the waiting-line saturation cliff. The correction removes
the mean per-leg travel bias entirely ($+0.744\to0.000$ time units)
and halves the end-to-end makespan drift between surrogate and
executor (mean absolute error $24.3\to12.1$ on held-out schedules).
The calibration is a frozen artifact: it is fitted once before
training, adds no learned components to the training loop, and is
shared by \emph{all} policies and by the dispatching-rule baseline of
Section~\ref{sec:training}, so every method plans in the same
corrected world. After training converges, the fit can be re-estimated
from the new policies' schedules to confirm stability.

% =========================================================================

\section{Policy Optimization and Refinement}
\label{sec:training}

\subsection{Objective and Gradient Estimator}

The policy $\pi_\theta$ is trained to minimize the expected executed
makespan $\mathbb{E}_{\pi_\theta}[C_{\max}(\Phi(\bm{a},\sigma))]$ with
the REINFORCE estimator \cite{williams1992reinforce}. During training,
episodes are generated by sampling from the masked policy
\eqref{eq:policy} on the calibrated surrogate; at validation and
deployment, actions are selected greedily. The gradient quality of
REINFORCE hinges on its baseline; we use a counterfactual baseline
built from a strong dispatching rule and evaluated by the executor
itself.

\subsection{Counterfactual Dispatch-Rule Baseline}
\label{sec:baseline}

Let $R$ denote a deterministic dispatching heuristic (job rule and AMR
rule; we use earliest-completion for both) and let
$R(\sigma_{1:t})$ denote the complete solution obtained by executing
the partial sequence $\sigma_{1:t}$ and completing all remaining
operations with $R$. Define the shorthand
$C[\cdot]=C_{\max}(\Phi(\cdot))$. The advantage of the action taken at
step $t$ is the improvement of committing that action and then
following the rule, relative to following the rule immediately:
\begin{equation}
A_t \;=\; C\big[R(\sigma_{1:t-1})\big]\;-\;
C\big[R(\sigma_{1:t-1}\oplus a_t)\big].
\label{eq:advantage}
\end{equation}
Both completions are evaluated by the collision-aware executor, so the
learning signal is grounded in true solution quality rather than in
surrogate clocks. Unlike episode-level rollout baselines
\cite{rennie2017self,kool2019attention}, \eqref{eq:advantage} assigns
credit per decision: an early pickup that congests a dock is penalized
at the step where it is committed, not diluted across the episode.
The price is $\mathcal{O}(n)$ executor evaluations per episode, which
dominates training cost; we consider the sharper credit assignment a
favorable trade at the instance sizes studied here, and
Section~\ref{sec:training-details} lists a cheaper actor--critic
alternative.

\subsection{Loss Function}

Advantages are standardized to zero mean and unit variance across all
steps of a batch of $B$ episodes. The per-batch loss combines the
policy gradient with a load-balancing shaping term and an entropy
bonus:
\begin{equation}
\mathcal{L}(\theta)=
-\frac{1}{B}\sum_{b=1}^{B}\sum_{t=1}^{2n}
\Big[\log\pi_\theta(a_t^b|s_t^b)\,
\big(\hat{A}_t^{\,b}+\lambda_{\mathrm{lb}}A_t^{\mathrm{lb},b}\big)
+\lambda_{H}\,\mathcal{H}_t^{b}\Big],
\label{eq:loss}
\end{equation}
where $\hat{A}_t$ is the standardized advantage,
$\mathcal{H}_t$ is the entropy of the masked policy over feasible
actions, and the shaping term
$A^{\mathrm{lb}}_t=\min_{k'}c_{k'}-c_{k}$ (for pickups; zero for
deliveries) compares the assignment count $c_k$ of the chosen AMR with
the least-loaded AMR, discouraging degenerate fleet utilization early
in training \cite{ng1999shaping}. We use
$\lambda_{\mathrm{lb}}=0.1$ and $\lambda_{H}=0.01$.

\begin{algorithm}[t]
\caption{Training with counterfactual dispatch baseline}
\label{alg:training}
\begin{algorithmic}[1]
\Require policy $\pi_\theta$, rule $R$, executor $\Phi$, surrogate $\hat{\Psi}$
\State fit calibration $(\alpha,\beta,\varphi)$ offline; freeze (Sec.~\ref{sec:calibration})
\For{$\mathrm{epoch}=1,\dots,E$}
  \For{$b=1,\dots,B$}
    \State sample instance; roll out $\sigma^b\!\sim\!\pi_\theta$ on $\hat{\Psi}$
    \For{$t=1,\dots,2n$}
      \State $A^b_t \gets C[R(\sigma^b_{1:t-1})]-C[R(\sigma^b_{1:t-1}\oplus a^b_t)]$
    \EndFor
  \EndFor
  \State standardize $\{A^b_t\}$ over the batch
  \State update $\theta$ by Adam on $\mathcal{L}(\theta)$ in \eqref{eq:loss}; clip $\lVert\nabla\rVert\le1$
  \If{$\mathrm{epoch}\bmod T_{\mathrm{val}}=0$}
    \State greedy-decode held-out instances; score by
    $C_{\max}+\kappa\cdot(\text{infeasible ops})$; keep best
  \EndIf
\EndFor
\end{algorithmic}
\end{algorithm}

\subsection{Training Protocol}
\label{sec:training-details}

Algorithm~\ref{alg:training} summarizes the procedure. Only the actor
parameters (encoders, GIN layers, operation embedding, and actor head)
receive gradients under REINFORCE; the critic \eqref{eq:critic} is
reserved for an optional proximal-policy-optimization variant
\cite{schulman2017ppo} that replaces \eqref{eq:advantage} with
critic-based advantages when the executor evaluations of the stepwise
baseline are prohibitive. We train with Adam \cite{kingma2015adam}
(learning rate $3\times10^{-4}$, batch $B=16$, gradient-norm clip
$1.0$) for $2{,}000$ epochs. Every $T_{\mathrm{val}}=50$ epochs the
policy is evaluated by greedy decoding on a fixed held-out validation
set, disjoint from the training pool, and scored by the executed
makespan plus a penalty of $\kappa=10^{3}$ per infeasible operation;
the checkpoint with the best validation score is retained. Computing
the entropy in \eqref{eq:loss} over the feasible-action support only
keeps the bonus comparable as the mask tightens toward the end of an
episode.

\subsection{Post-Processing Local Search}
\label{sec:localsearch}

At deployment, the greedy-decoded schedule is refined by a two-stage
first-improvement local search that reuses the neighborhood operators
of our genetic-algorithm baseline, so that quality differences between
methods are attributable to the constructed solutions rather than to
the refinement. A move either swaps two operations in $\sigma$, at
least one of which belongs to the AMR that currently determines the
makespan, or relocates a contiguous block of same-AMR, same-type
operations; the AMR assignment $\bm{a}$ is left unchanged. The first
stage performs $I_{\mathrm{cf}}=10^{3}$ trials under the fast
collision-free evaluator to explore broadly; the second performs
$I_{\mathrm{ca}}=10^{2}$ trials under the full executor $\Phi$ to
polish under true routing. The refinement adds approximately two
seconds per instance and is entirely optional: all learned-policy
results are reported both with and without it.

% =========================================================================
% Suggested BibTeX entries (replace keys above with your bibliography):
%
% @inproceedings{xu2019gin, Xu, Hu, Leskovec, Jegelka, "How Powerful Are
%   Graph Neural Networks?", ICLR 2019}
% @inproceedings{zhang2020l2d, Zhang et al., "Learning to Dispatch for
%   Job Shop Scheduling via Deep Reinforcement Learning", NeurIPS 2020}
% @article{song2023fjsp, Song et al., "Flexible Job-Shop Scheduling via
%   Graph Neural Network and Deep Reinforcement Learning", IEEE TII 2023}
% @inproceedings{li2018deeper, Li, Han, Wu, "Deeper Insights into Graph
%   Convolutional Networks for Semi-Supervised Learning", AAAI 2018}
% @inproceedings{oono2020oversmoothing, Oono, Suzuki, "Graph Neural
%   Networks Exponentially Lose Expressive Power for Node
%   Classification", ICLR 2020}
% @inproceedings{stern2019mapf, Stern et al., "Multi-Agent Pathfinding:
%   Definitions, Variants, and Benchmarks", SoCS 2019}
% @article{williams1992reinforce, Williams, "Simple Statistical
%   Gradient-Following Algorithms for Connectionist Reinforcement
%   Learning", Machine Learning 8, 1992}
% @inproceedings{rennie2017self, Rennie et al., "Self-Critical Sequence
%   Training for Image Captioning", CVPR 2017}
% @inproceedings{kool2019attention, Kool, van Hoof, Welling, "Attention,
%   Learn to Solve Routing Problems!", ICLR 2019}
% @inproceedings{ng1999shaping, Ng, Harada, Russell, "Policy Invariance
%   Under Reward Transformations: Theory and Application to Reward
%   Shaping", ICML 1999}
% @article{schulman2017ppo, Schulman et al., "Proximal Policy
%   Optimization Algorithms", arXiv:1707.06347}
% @inproceedings{kingma2015adam, Kingma, Ba, "Adam: A Method for
%   Stochastic Optimization", ICLR 2015}
% =========================================================================
 inside an IEEEtran doc.
% Requires in the preamble:
%   \usepackage{amsmath,amssymb,bm}
%   \usepackage{booktabs,array,multirow,graphicx}
%   \usepackage{algorithm}
%   \usepackage{algpseudocode}
% Citation keys are placeholders -- suggested BibTeX entries are listed in
% comments at the end. "CD-GNN" (Cross-Docking GNN) is a placeholder model
% name; rename globally if you prefer.
% Sections: Problem Formulation, Architecture, Surrogate Calibration,
% Policy Optimization (incl. post-processing local search).
% =========================================================================

\section{Problem Formulation}
\label{sec:problem}

\begin{table*}[t]
\caption{Summary of Notation}
\label{tab:notation}
\centering
\footnotesize
\begin{tabular}{@{}l p{0.295\textwidth} @{\qquad} l p{0.295\textwidth}@{}}
\toprule
Symbol & Meaning & Symbol & Meaning \\
\midrule
\multicolumn{2}{@{}l}{\itshape Problem setting} &
\multicolumn{2}{@{}l}{\itshape Architecture} \\
$\mathcal{W}$ & grid workspace (four-connected unit cells) &
  $\bm{x}^{A}_k,\bm{x}^{J}_j$ & AMR / job input features ($\mathbb{R}^{8},\mathbb{R}^{16}$) \\
$\Delta(u,v)$ & Manhattan distance between positions $u,v$ &
  $h$ & hidden width of all embeddings ($h{=}128$) \\
$\mathcal{J},\ n$ & job set; number of jobs &
  $\bm{h}^{A}_k$ & AMR embedding \\
$\mathcal{M},\ m$ & AMR fleet; fleet size &
  $\bm{g}^{(l)}_j$ & job embedding after $l$ GIN layers \\
$\mathcal{T}$ & set of material types &
  $\mathcal{G}_t=(\mathcal{J},E_t)$ & dynamic precedence graph at step $t$ \\
$\mathcal{D}^{\mathrm{in}},\mathcal{D}^{\mathrm{out}},\mathcal{D}$ & inbound docks; outbound docks; their union &
  $\mathcal{P}(j)$ & committed predecessors of job $j$ in $\mathcal{G}_t$ \\
$\tau_j$ & material type of job $j$ &
  $\epsilon^{(l)},\ L$ & learnable GIN self-weight; number of GIN layers \\
$o_j,\ d_j$ & inbound and outbound dock of job $j$ &
  $\mathrm{LN}$ & layer normalization \\
$r_j$ & release time of job $j$ &
  $\bm{e}_{\omega}$ & learned operation-type embedding \\
$p_j=p(\tau_j)$ & dock service duration of job $j$ &
  $z_{\omega,k,j}$ & logit of action $(\omega,k,j)$ \\
$\pi_j,\ \delta_j$ & pickup and delivery operations of job $j$ &
  $V_{\phi}$ & critic (state-value) network, parameters $\phi$ \\
$Q$ & per-type AMR carrying capacity &
  & \\
$\bm{a},\ a_j$ & AMR assignment vector; AMR assigned to job $j$ &
  \multicolumn{2}{@{}l}{\itshape Training and refinement} \\
$\sigma,\ \sigma_t$ & operation sequence; prefix committed up to step $t$ &
  $R,\ R(\sigma_{1:t})$ & dispatching rule; completion of prefix $\sigma_{1:t}$ by $R$ \\
$s^{\pi}_j,c^{\pi}_j$ & start / completion of pickup service of job $j$ &
  $C[\cdot]$ & shorthand for $C_{\max}(\Phi(\cdot))$ \\
$s^{\delta}_j,c^{\delta}_j$ & start / completion of delivery service of job $j$ &
  $A_t,\ \hat{A}_t$ & stepwise counterfactual advantage; standardized \\
$t_k(u,v)$ & realized travel time of AMR $k$ from $u$ to $v$ &
  $A^{\mathrm{lb}}_t$ & load-balance shaping advantage \\
$\Phi$ & collision-aware executor &
  $\lambda_{\mathrm{lb}},\ \lambda_{H}$ & shaping and entropy coefficients \\
$C_{\max},\ F_k$ & makespan; finish time of AMR $k$ incl.\ depot return &
  $\mathcal{H}_t$ & entropy of the masked policy at step $t$ \\
$W_d$ & number of physical waiting cells at dock $d$ &
  $\mathcal{L}(\theta)$ & training loss \\
\multicolumn{2}{@{}l}{\itshape Sequential decision process} &
  $B,\ E$ & batch size (episodes); number of training epochs \\
$s_t,\ a_t$ & state and action at decision epoch $t$ &
  $T_{\mathrm{val}},\ \kappa$ & validation interval; infeasibility penalty \\
$\omega$ & operation type, $\omega\in\{\textsc{pick},\textsc{deliver}\}$ &
  $I_{\mathrm{cf}},\ I_{\mathrm{ca}}$ & collision-free / collision-aware refinement budgets \\
$\Lambda(s_t)$ & feasible-action set (mask) in state $s_t$ &
  $\mathbf{1}[\cdot]$ & indicator function \\
$\hat{\Psi}$ & calibrated surrogate transition model &
  & \\
$\mathcal{E}_d$ & committed service events $(s_e,c_e,j_e,\omega_e)$ at dock $d$ &
  & \\
$\hat{b}_k,\ x_k$ & next-available time and position of AMR $k$ &
  & \\
$\rho$ & ready time of an AMR at a dock &
  & \\
$\hat{t}(u,v)$ & calibrated travel-time estimate \eqref{eq:travelcal} &
  & \\
$\hat{w}_d(\rho)$ & queue-aware dock-wait estimate \eqref{eq:waitcal} &
  & \\
$q_d(\rho)$ & services still unfinished at time $\rho$ at dock $d$ &
  & \\
$\alpha,\ \beta$ & travel-calibration scale and offset &
  & \\
$\varphi(\cdot)$ & dock-wait penalty table over queue depths &
  & \\
$\pi_{\theta}$ & scheduling policy with parameters $\theta$ &
  & \\
\bottomrule
\end{tabular}
\end{table*}

\subsection{System Description}

We consider the internal transfer problem of a cross-docking terminal
served by a fleet of autonomous mobile robots (AMRs). Table~\ref{tab:notation}
summarizes the notation used throughout the paper. The terminal floor
is modeled as a four-connected unit grid $\mathcal{W}$ in which a set of
\emph{inbound docks} $\mathcal{D}^{\mathrm{in}}$ is located along one
boundary and a set of \emph{outbound docks} (processing stations)
$\mathcal{D}^{\mathrm{out}}$ along the opposite boundary; the AMR depots
lie between them. Consistent with cross-docking practice, the floor is
free of static obstacles, so congestion arises from vehicle interference
and dock contention rather than from the layout itself. AMRs travel at
unit speed between adjacent cells, and the distance between two
positions $u,v$ is the Manhattan metric $\Delta(u,v)$.

A set of transfer jobs $\mathcal{J}=\{1,\dots,n\}$ must be moved across
the terminal by a fleet $\mathcal{M}=\{1,\dots,m\}$ of homogeneous AMRs.
Each job $j\in\mathcal{J}$ represents one unit load of a material type
$\tau_j\in\mathcal{T}$ that becomes available at its inbound dock
$o_j\in\mathcal{D}^{\mathrm{in}}$ at release time $r_j\ge 0$ and must be
delivered to a designated outbound dock
$d_j\in\mathcal{D}^{\mathrm{out}}$. Handling a unit load of type $\tau$
occupies a dock for a type-dependent service duration $p(\tau)$; we
write $p_j=p(\tau_j)$ and assume the loading service at the inbound dock
and the unloading/processing service at the outbound dock take the same
time, reflecting symmetric handling of a unit load. Each job therefore
consists of an ordered pair of operations: a \emph{pickup} operation
$\pi_j$ executed at $o_j$ and a \emph{delivery} operation $\delta_j$
executed at $d_j$, both performed by the same AMR (no transshipment).
An AMR may consolidate loads subject to a per-type carrying capacity
$Q$: at no time may it carry more than $Q$ units of any single material
type.

Docks are exclusive single-server resources: at most one AMR may be in
service at a dock at any time, and services are non-preemptive. An AMR
arriving at an occupied dock queues in a bounded set of waiting cells
adjacent to that dock; when the waiting line is saturated, the AMR must
hold at an upstream position and re-approach later. Finally, AMR motion
must be collision-free: at any unit time step, each grid cell may be
occupied by at most one AMR.

\subsection{Formal Statement}

A solution is a pair $(\bm{a},\sigma)$, where
$\bm{a}=(a_1,\dots,a_n)\in\mathcal{M}^{n}$ assigns each job to an AMR
and $\sigma$ is a total order over the $2n$ operations
$\{\pi_j,\delta_j\}_{j\in\mathcal{J}}$ that determines the sequence in
which operations are committed. Let $s^{\pi}_j,c^{\pi}_j$
(resp.\ $s^{\delta}_j,c^{\delta}_j$) denote the service start and
completion times of $\pi_j$ (resp.\ $\delta_j$) induced by executing
$(\bm{a},\sigma)$. A feasible execution satisfies:
\begin{align}
c^{\pi}_j &= s^{\pi}_j + p_j, \qquad
c^{\delta}_j = s^{\delta}_j + p_j,
&& \forall j\in\mathcal{J}, \label{eq:nonpreempt}\\
s^{\pi}_j &\ge r_j,
&& \forall j\in\mathcal{J}, \label{eq:release}\\
s^{\delta}_j &\ge c^{\pi}_j,
&& \forall j\in\mathcal{J}, \label{eq:precedence}
\end{align}
\begin{align}
\big|\{\,j : a_j{=}k,\ \tau_j{=}\tau,\ c^{\pi}_j \le u < c^{\delta}_j \}\big| &\le Q,
&& \forall k,\tau,u, \label{eq:capacity}\\
\big[s_e, c_e\big) \cap \big[s_{e'}, c_{e'}\big) &= \emptyset,
&& \forall e\ne e' \ \text{at the same dock}, \label{eq:exclusive}
\end{align}
where \eqref{eq:nonpreempt} enforces non-preemptive services,
\eqref{eq:release} release dates, \eqref{eq:precedence} pickup-before-%
delivery precedence, \eqref{eq:capacity} the per-type carrying capacity,
and \eqref{eq:exclusive} dock exclusivity over all service events $e$.
In addition, for any two operations $u\prec v$ executed consecutively by
the same AMR $k$,
\begin{equation}
s_v \;\ge\; c_u + t_k\big(\mathrm{loc}(u),\mathrm{loc}(v)\big),
\label{eq:travel}
\end{equation}
where $t_k(\cdot,\cdot)\ge\Delta(\cdot,\cdot)$ is the realized travel
time of AMR $k$, which may exceed the Manhattan distance due to
collision avoidance, queuing detours, and dock-clearing maneuvers.

The realized travel and waiting times are not fixed constants but are
determined by a high-fidelity \emph{executor} $\Phi$, which routes every
AMR leg with space--time $\mathrm{A}^{*}$ search over a shared
reservation table (at most one AMR per cell per time step), manages the
bounded dock waiting lines, and inserts upstream holding maneuvers when
lines are saturated. The problem is therefore a simulation-based
optimization: find
\begin{equation}
(\bm{a}^{*},\sigma^{*}) \;=\;
\arg\min_{(\bm{a},\sigma)} \; C_{\max}\big(\Phi(\bm{a},\sigma)\big),
\label{eq:objective}
\end{equation}
where the makespan $C_{\max}=\max_{k\in\mathcal{M}} F_k$ and $F_k$ is
the time at which AMR $k$ completes its final activity, including its
return trip to the depot. The problem generalizes both the capacitated
pickup-and-delivery problem and parallel-machine scheduling with
sequence-dependent setups, and is therefore NP-hard; the coupling to
collision-free multi-agent routing further compounds its difficulty
\cite{stern2019mapf}.

\subsection{Sequential Decision Formulation}
\label{sec:mdp}

We cast the construction of $(\bm{a},\sigma)$ as an episodic sequential
decision process with exactly $2n$ decision epochs, one per committed
operation. At epoch $t$, the state
\begin{equation}
s_t=\big(\{\bm{x}^{A}_{k}\},\, \{\bm{x}^{J}_{j}\},\,
\{\mathcal{E}_d\},\, \sigma_{t}\big)
\label{eq:state}
\end{equation}
collects the AMR states (position, next-available time, per-type
inventory), the job states (unpicked / onboard / delivered, carrier
identity), the set of committed dock service events
$\mathcal{E}_d=\{(s_e,c_e,j_e,\omega_e)\}$ at each dock $d$, and the
partial operation sequence $\sigma_t$.

An action is a triple
\begin{equation}
a_t=(\omega,k,j)\in\{\textsc{pick},\textsc{deliver}\}\times\mathcal{M}\times\mathcal{J},
\label{eq:action}
\end{equation}
giving a flat action space of size $2mn$ restricted by a feasibility
mask $\Lambda(s_t)$: $(\textsc{pick},k,j)$ is legal iff $j$ is unpicked
and AMR $k$ carries fewer than $Q$ units of type $\tau_j$;
$(\textsc{deliver},k,j)$ is legal iff $j$ is onboard AMR $k$. The mask
guarantees by construction that every complete episode yields a feasible
solution, so no repair mechanism is required.

Transitions follow a deterministic surrogate model $\hat{\Psi}$ that
advances the acting AMR's clock using a \emph{calibrated} travel-time
estimate $\hat{t}(\cdot,\cdot)$ and a queue-aware dock-wait estimate
$\hat{w}_d(\cdot)$ fitted offline against the executor $\Phi$
(Section~\ref{sec:calibration}); committing an action appends the
corresponding service window to $\mathcal{E}_d$. For a pickup, the
service start of job $j$ by AMR $k$ located at $x_k$ with next-available
time $\hat{b}_k$ is
\begin{equation}
s^{\pi}_j=\underbrace{\max\!\big(\hat{b}_k+\hat{t}(x_k,o_j),\, r_j\big)}_{\text{ready time }\rho}
\;+\;\hat{w}_{o_j}(\rho),
\label{eq:surrogate}
\end{equation}
and analogously for deliveries at $d_j$. After $2n$ steps the completed
solution is evaluated by the executor, and the learning objective is to
minimize the expected executed makespan
$\mathbb{E}_{\pi_\theta}\!\left[C_{\max}(\Phi(\bm{a},\sigma))\right]$;
the policy-gradient training procedure and its dispatch-rule baseline
are described in Section~\ref{sec:training}.

% =========================================================================

\section{Sparse-Graph Multi-Pointer Dispatcher}
\label{sec:architecture}

\subsection{Overview and Design Rationale}

Learning-to-dispatch architectures for shop scheduling encode
operations as nodes of a disjunctive graph and learn a construction
policy over it \cite{zhang2020l2d,song2023fjsp}. Transferring this
paradigm to AMR-served cross-docking raises a representation question:
how should the tight resource coupling---$m$ vehicles contending for
$|\mathcal{D}|$ single-server docks under collision-constrained
motion---enter the graph? A seemingly natural answer is to join every
pair of jobs that share a dock. We found this to be
counterproductive: with $n/|\mathcal{D}|$ jobs per dock, such
contention edges form dense cliques, and message passing over them
drives all job embeddings toward a common average (pairwise cosine
similarity above $0.99$ at initialization in our measurements),
leaving the policy unable to discriminate between candidate jobs
precisely at the opening decisions that dominate the makespan.
Over-smoothing under dense neighborhoods is a known failure mode of
deep graph networks \cite{li2018deeper,oono2020oversmoothing}; in this
problem it appears already at three message-passing rounds.

The proposed scheduler, hereafter \textsc{CD-GNN}, therefore keeps the
graph \emph{sparse and dynamic}: edges encode only scheduling
precedences that have actually been committed, while resource
congestion enters through calibrated, decision-grade node features.
The model comprises (i)~a job encoder over a dynamic
\emph{precedence graph} processed by residual graph isomorphism
network (GIN) layers \cite{xu2019gin}; (ii)~a lightweight fleet
encoder; and (iii)~a factored multi-pointer actor that scores each
masked (operation, AMR, job) triple \cite{kool2019attention}. The full
model is deliberately lean: $2.8\times10^{5}$ parameters
($2.5\times10^{5}$ actor, $3.3\times10^{4}$ critic) with hidden width
$h=128$ throughout, permutation-invariant in the job and AMR sets, and
applicable unchanged to any instance size.

% TODO: architecture overview figure
% \begin{figure*}[t]
%   \centering
%   \includegraphics[width=\textwidth]{figures/cdgnn_architecture}
%   \caption{Overview of \textsc{CD-GNN}. Job features (including
%   calibrated congestion estimates) are encoded and refined over the
%   dynamic precedence graph by residual GIN layers; AMR features are
%   encoded independently; the multi-pointer actor scores each feasible
%   (operation, AMR, job) triple under the feasibility mask.}
%   \label{fig:architecture}
% \end{figure*}

\subsection{Congestion-Aware State Representation}
\label{sec:features}

At every decision epoch the state \eqref{eq:state} is rendered into two
feature families, summarized in Table~\ref{tab:features}. All features
are normalized to comparable magnitudes by fixed scales.

\emph{Job nodes} ($\bm{x}^{J}_j\in\mathbb{R}^{16}$). Each job encodes
its service duration, lifecycle status (unpicked/onboard/delivered) and
carrier identity, its material type, and the coordinates of both its
supply position and its destination station, giving the network direct
geometric grounding. Resource congestion enters through four derived
features: the current availability delays of the job's inbound and
outbound docks, and the number of still-unpicked competitor jobs
committed to each of the two docks. A calibrated completion-time
estimate (Section~\ref{sec:calibration}) summarizes how quickly the job
could finish if served next, providing a value-like signal that the
network would otherwise need several layers to reconstruct.

\emph{AMR nodes} ($\bm{x}^{A}_k\in\mathbb{R}^{8}$). Each AMR encodes a
busy indicator and its time until available, per-type inventory
ratios, its coordinates, and the number of jobs currently assigned to
it, which exposes fleet-utilization imbalance to the policy.

\begin{table}[t]
\caption{Node Features of \textsc{CD-GNN}}
\label{tab:features}
\centering
\footnotesize
\begin{tabular}{@{}llc@{}}
\toprule
& Feature & Scale \\
\midrule
\multirow{12}{*}{\rotatebox{90}{Job $\bm{x}^{J}_j\!\in\!\mathbb{R}^{16}$}}
& existence flag & -- \\
& service duration $p_j$ & $1/25$ \\
& carrier index (0 if unpicked) & $1/m$ \\
& destination position $(x,y)$ ($\times2$) & $1/10$ \\
& supply position $(x,y)$ ($\times2$) & $1/10$ \\
& material type one-hot ($\times3$) & -- \\
& lifecycle status $\{0,\tfrac12,1\}$ & -- \\
& completion-time estimate & $1/500$ \\
& inbound dock availability delay & $1/100$ \\
& outbound dock availability delay & $1/100$ \\
& inbound dock committed queue & $1/m$ \\
& outbound dock committed queue & $1/m$ \\
\midrule
\multirow{7}{*}{\rotatebox{90}{AMR $\bm{x}^{A}_k\!\in\!\mathbb{R}^{8}$}}
& busy indicator & -- \\
& time until available & $1/50$ \\
& inventory ratio, per type ($\times3$) & $1/Q$ \\
& position $(x,y)$ ($\times2$) & $1/10$ \\
& assigned-job count & $1/10$ \\
\bottomrule
\end{tabular}
\end{table}

\subsection{Dynamic Precedence Graph and Job Encoder}
\label{sec:graph}

Job features pass through a linear map
$\bm{g}^{(0)}_j=\bm{W}_{\!J}\,\bm{x}^{J}_j\in\mathbb{R}^{h}$ and AMR
features through a linear encoder
$\bm{h}^{A}_k=\bm{W}_{\!A}\,\bm{x}^{A}_k\in\mathbb{R}^{h}$.

Jobs are coupled by the \emph{committed} structure of the partial
schedule. Following the adding-arc construction of disjunctive-graph
dispatching \cite{zhang2020l2d}, the graph
$\mathcal{G}_t=(\mathcal{J},E_t)$ starts empty and grows one or two
directed arcs per decision: when an operation of job $j$ is committed,
an arc is added from the previously scheduled job on the same AMR to
$j$, and from the previously scheduled job on the same outbound dock to
$j$. Each job thus aggregates information from its immediate
predecessors $\mathcal{P}(j)$ on both resources it occupies---the
vehicle sequence and the dock sequence---while uncommitted jobs
interact only through the shared congestion features of
Section~\ref{sec:features}, never through dense same-dock edges.

Message passing applies $L=3$ residual GIN layers with layer
normalization,
\begin{equation}
\bm{g}^{(l+1)}_j=\bm{g}^{(l)}_j+
\mathrm{LN}\!\left(\mathrm{MLP}^{(l)}\!\Big(
(1{+}\epsilon^{(l)})\,\bm{g}^{(l)}_j+
\sum_{i\in\mathcal{P}(j)}\!\bm{g}^{(l)}_i
\Big)\right)\!,
\label{eq:gin}
\end{equation}
retaining the injective sum aggregation of GIN \cite{xu2019gin}: with
at most two predecessors per node, message magnitudes stay bounded
irrespective of instance size, so the scale-invariance concerns that
arise with dense contention edges do not apply.

\subsection{Multi-Pointer Action Scoring and Value Estimation}
\label{sec:actor}

Let $\bm{e}_{\omega}\in\mathbb{R}^{h}$ be a learned embedding of the
operation type. The actor scores every candidate triple by
\begin{equation}
z_{\omega,k,j}=\mathrm{MLP}_{\pi}\!\big(\big[\,\bm{e}_{\omega}\,\|\,
\bm{g}^{(L)}_{j}\,\|\,\bm{h}^{A}_{k}\big]\big),
\label{eq:logit}
\end{equation}
and the policy is the masked softmax
\begin{equation}
\pi_{\theta}(a\,|\,s_t)=
\frac{\exp(z_{a})\,\mathbf{1}[a\in\Lambda(s_t)]}
{\sum_{a'\in\Lambda(s_t)}\exp(z_{a'})}.
\label{eq:policy}
\end{equation}
The critic estimates the state value from pooled entity summaries,
\begin{equation}
V_{\phi}(s_t)=\mathrm{MLP}_{V}\!\big(\big[\,
\overline{\bm{g}}^{\mathrm{(unf)}} \,\|\, \overline{\bm{h}}^{A}\big]\big),
\label{eq:critic}
\end{equation}
where $\overline{\bm{g}}^{\mathrm{(unf)}}$ averages the embeddings of
unfinished jobs only. A full episode requires exactly $2n$ forward
passes; each pass scores all $2mn$ actions in a single batched
evaluation of \eqref{eq:logit}.

% =========================================================================

\section{Surrogate Model Calibration}
\label{sec:calibration}

\subsection{The Two-Simulator Gap}

The policy constructs schedules on the surrogate $\hat{\Psi}$ of
Section~\ref{sec:mdp}, whereas solution quality is defined by the
executor $\Phi$ of \eqref{eq:objective}. An uncalibrated surrogate that
advances clocks by Manhattan distances and single-scalar dock
availabilities is systematically optimistic: it ignores collision
detours, dock-clearing maneuvers, waiting-line approach paths, and the
upstream holds triggered when waiting lines saturate. The congestion
features of Section~\ref{sec:features}---and therefore every decision
the policy makes---would then be conditioned on clocks that drift
further from reality as congestion grows, precisely in the regime where
scheduling decisions matter most. Rather than paying the computational
cost of embedding the full executor in the training loop, we close most
of this gap with two lightweight corrections fitted offline.

\subsection{Calibrated Travel and Dock-Wait Estimators}

\emph{Travel.} Realized leg times are modeled as an affine function of
the Manhattan distance, clamped from below by the metric itself:
\begin{equation}
\hat{t}(u,v)=\max\!\big(\Delta(u,v),\ \alpha\,\Delta(u,v)+\beta\big).
\label{eq:travelcal}
\end{equation}

\emph{Dock waiting.} The wait experienced by an AMR that becomes ready
at a dock $d$ at time $\rho$ is modeled as the committed availability
delay plus a penalty indexed by the queue depth at readiness,
\begin{equation}
\hat{w}_d(\rho)=\max\!\big(0,\ \mathrm{avail}_d-\rho\big)
+\varphi\!\big(\min(q_d(\rho),\,3)\big),
\label{eq:waitcal}
\end{equation}
where $q_d(\rho)=|\{e\in\mathcal{E}_d : c_e>\rho\}|$ counts committed
services unfinished at $\rho$, and
$\varphi:\{0,1,2,3^{+}\}\to\mathbb{R}_{\ge0}$ is a monotone lookup
table. The bounded table shape is deliberate: waiting-line detours and
upstream holds produce a cost cliff once the line saturates, a
nonlinearity that a linear congestion model cannot express.

\subsection{Fitting Procedure}

Calibration data are obtained by replaying complete schedules through
both simulators and aligning them per operation. Schedules are drawn
from heterogeneous sources---five dispatching-rule pairs and the
current learned policies---across instance sizes
$n\in\{10,20,40,60\}$, so that the fitted correction covers the state
distribution that policies actually visit. For every operation, the
raw surrogate supplies the Manhattan leg distance, the base wait, and
the queue depth at readiness; the executor's logged timeline is
partitioned, per AMR, into legs terminating at each service start,
within which movement entries sum to the realized travel and the
remainder is attributed to waiting. Executions containing infeasible
operations are discarded. The travel parameters $(\alpha,\beta)$ are
fitted by least squares and the penalties $\varphi(\cdot)$ as clipped
bin means with monotonicity enforced across depths.

In our environment, fitting on $10{,}360$ operation pairs from $159$
schedules yielded $\alpha=0.963$, $\beta=1.12$ ($R^{2}=0.91$) and
$\varphi=(0.04,\,0.04,\,1.66,\,3.57)$; the step between depths one and
two locates the waiting-line saturation cliff. The correction removes
the mean per-leg travel bias entirely ($+0.744\to0.000$ time units)
and halves the end-to-end makespan drift between surrogate and
executor (mean absolute error $24.3\to12.1$ on held-out schedules).
The calibration is a frozen artifact: it is fitted once before
training, adds no learned components to the training loop, and is
shared by \emph{all} policies and by the dispatching-rule baseline of
Section~\ref{sec:training}, so every method plans in the same
corrected world. After training converges, the fit can be re-estimated
from the new policies' schedules to confirm stability.

% =========================================================================

\section{Policy Optimization and Refinement}
\label{sec:training}

\subsection{Objective and Gradient Estimator}

The policy $\pi_\theta$ is trained to minimize the expected executed
makespan $\mathbb{E}_{\pi_\theta}[C_{\max}(\Phi(\bm{a},\sigma))]$ with
the REINFORCE estimator \cite{williams1992reinforce}. During training,
episodes are generated by sampling from the masked policy
\eqref{eq:policy} on the calibrated surrogate; at validation and
deployment, actions are selected greedily. The gradient quality of
REINFORCE hinges on its baseline; we use a counterfactual baseline
built from a strong dispatching rule and evaluated by the executor
itself.

\subsection{Counterfactual Dispatch-Rule Baseline}
\label{sec:baseline}

Let $R$ denote a deterministic dispatching heuristic (job rule and AMR
rule; we use earliest-completion for both) and let
$R(\sigma_{1:t})$ denote the complete solution obtained by executing
the partial sequence $\sigma_{1:t}$ and completing all remaining
operations with $R$. Define the shorthand
$C[\cdot]=C_{\max}(\Phi(\cdot))$. The advantage of the action taken at
step $t$ is the improvement of committing that action and then
following the rule, relative to following the rule immediately:
\begin{equation}
A_t \;=\; C\big[R(\sigma_{1:t-1})\big]\;-\;
C\big[R(\sigma_{1:t-1}\oplus a_t)\big].
\label{eq:advantage}
\end{equation}
Both completions are evaluated by the collision-aware executor, so the
learning signal is grounded in true solution quality rather than in
surrogate clocks. Unlike episode-level rollout baselines
\cite{rennie2017self,kool2019attention}, \eqref{eq:advantage} assigns
credit per decision: an early pickup that congests a dock is penalized
at the step where it is committed, not diluted across the episode.
The price is $\mathcal{O}(n)$ executor evaluations per episode, which
dominates training cost; we consider the sharper credit assignment a
favorable trade at the instance sizes studied here, and
Section~\ref{sec:training-details} lists a cheaper actor--critic
alternative.

\subsection{Loss Function}

Advantages are standardized to zero mean and unit variance across all
steps of a batch of $B$ episodes. The per-batch loss combines the
policy gradient with a load-balancing shaping term and an entropy
bonus:
\begin{equation}
\mathcal{L}(\theta)=
-\frac{1}{B}\sum_{b=1}^{B}\sum_{t=1}^{2n}
\Big[\log\pi_\theta(a_t^b|s_t^b)\,
\big(\hat{A}_t^{\,b}+\lambda_{\mathrm{lb}}A_t^{\mathrm{lb},b}\big)
+\lambda_{H}\,\mathcal{H}_t^{b}\Big],
\label{eq:loss}
\end{equation}
where $\hat{A}_t$ is the standardized advantage,
$\mathcal{H}_t$ is the entropy of the masked policy over feasible
actions, and the shaping term
$A^{\mathrm{lb}}_t=\min_{k'}c_{k'}-c_{k}$ (for pickups; zero for
deliveries) compares the assignment count $c_k$ of the chosen AMR with
the least-loaded AMR, discouraging degenerate fleet utilization early
in training \cite{ng1999shaping}. We use
$\lambda_{\mathrm{lb}}=0.1$ and $\lambda_{H}=0.01$.

\begin{algorithm}[t]
\caption{Training with counterfactual dispatch baseline}
\label{alg:training}
\begin{algorithmic}[1]
\Require policy $\pi_\theta$, rule $R$, executor $\Phi$, surrogate $\hat{\Psi}$
\State fit calibration $(\alpha,\beta,\varphi)$ offline; freeze (Sec.~\ref{sec:calibration})
\For{$\mathrm{epoch}=1,\dots,E$}
  \For{$b=1,\dots,B$}
    \State sample instance; roll out $\sigma^b\!\sim\!\pi_\theta$ on $\hat{\Psi}$
    \For{$t=1,\dots,2n$}
      \State $A^b_t \gets C[R(\sigma^b_{1:t-1})]-C[R(\sigma^b_{1:t-1}\oplus a^b_t)]$
    \EndFor
  \EndFor
  \State standardize $\{A^b_t\}$ over the batch
  \State update $\theta$ by Adam on $\mathcal{L}(\theta)$ in \eqref{eq:loss}; clip $\lVert\nabla\rVert\le1$
  \If{$\mathrm{epoch}\bmod T_{\mathrm{val}}=0$}
    \State greedy-decode held-out instances; score by
    $C_{\max}+\kappa\cdot(\text{infeasible ops})$; keep best
  \EndIf
\EndFor
\end{algorithmic}
\end{algorithm}

\subsection{Training Protocol}
\label{sec:training-details}

Algorithm~\ref{alg:training} summarizes the procedure. Only the actor
parameters (encoders, GIN layers, operation embedding, and actor head)
receive gradients under REINFORCE; the critic \eqref{eq:critic} is
reserved for an optional proximal-policy-optimization variant
\cite{schulman2017ppo} that replaces \eqref{eq:advantage} with
critic-based advantages when the executor evaluations of the stepwise
baseline are prohibitive. We train with Adam \cite{kingma2015adam}
(learning rate $3\times10^{-4}$, batch $B=16$, gradient-norm clip
$1.0$) for $2{,}000$ epochs. Every $T_{\mathrm{val}}=50$ epochs the
policy is evaluated by greedy decoding on a fixed held-out validation
set, disjoint from the training pool, and scored by the executed
makespan plus a penalty of $\kappa=10^{3}$ per infeasible operation;
the checkpoint with the best validation score is retained. Computing
the entropy in \eqref{eq:loss} over the feasible-action support only
keeps the bonus comparable as the mask tightens toward the end of an
episode.

\subsection{Post-Processing Local Search}
\label{sec:localsearch}

At deployment, the greedy-decoded schedule is refined by a two-stage
first-improvement local search that reuses the neighborhood operators
of our genetic-algorithm baseline, so that quality differences between
methods are attributable to the constructed solutions rather than to
the refinement. A move either swaps two operations in $\sigma$, at
least one of which belongs to the AMR that currently determines the
makespan, or relocates a contiguous block of same-AMR, same-type
operations; the AMR assignment $\bm{a}$ is left unchanged. The first
stage performs $I_{\mathrm{cf}}=10^{3}$ trials under the fast
collision-free evaluator to explore broadly; the second performs
$I_{\mathrm{ca}}=10^{2}$ trials under the full executor $\Phi$ to
polish under true routing. The refinement adds approximately two
seconds per instance and is entirely optional: all learned-policy
results are reported both with and without it.

% =========================================================================
% Suggested BibTeX entries (replace keys above with your bibliography):
%
% @inproceedings{xu2019gin, Xu, Hu, Leskovec, Jegelka, "How Powerful Are
%   Graph Neural Networks?", ICLR 2019}
% @inproceedings{zhang2020l2d, Zhang et al., "Learning to Dispatch for
%   Job Shop Scheduling via Deep Reinforcement Learning", NeurIPS 2020}
% @article{song2023fjsp, Song et al., "Flexible Job-Shop Scheduling via
%   Graph Neural Network and Deep Reinforcement Learning", IEEE TII 2023}
% @inproceedings{li2018deeper, Li, Han, Wu, "Deeper Insights into Graph
%   Convolutional Networks for Semi-Supervised Learning", AAAI 2018}
% @inproceedings{oono2020oversmoothing, Oono, Suzuki, "Graph Neural
%   Networks Exponentially Lose Expressive Power for Node
%   Classification", ICLR 2020}
% @inproceedings{stern2019mapf, Stern et al., "Multi-Agent Pathfinding:
%   Definitions, Variants, and Benchmarks", SoCS 2019}
% @article{williams1992reinforce, Williams, "Simple Statistical
%   Gradient-Following Algorithms for Connectionist Reinforcement
%   Learning", Machine Learning 8, 1992}
% @inproceedings{rennie2017self, Rennie et al., "Self-Critical Sequence
%   Training for Image Captioning", CVPR 2017}
% @inproceedings{kool2019attention, Kool, van Hoof, Welling, "Attention,
%   Learn to Solve Routing Problems!", ICLR 2019}
% @inproceedings{ng1999shaping, Ng, Harada, Russell, "Policy Invariance
%   Under Reward Transformations: Theory and Application to Reward
%   Shaping", ICML 1999}
% @article{schulman2017ppo, Schulman et al., "Proximal Policy
%   Optimization Algorithms", arXiv:1707.06347}
% @inproceedings{kingma2015adam, Kingma, Ba, "Adam: A Method for
%   Stochastic Optimization", ICLR 2015}
% =========================================================================
 inside an IEEEtran doc.
% Requires in the preamble:
%   \usepackage{amsmath,amssymb,bm}
%   \usepackage{booktabs,array,multirow,graphicx}
%   \usepackage{algorithm}
%   \usepackage{algpseudocode}
% Citation keys are placeholders -- suggested BibTeX entries are listed in
% comments at the end. "CD-GNN" (Cross-Docking GNN) is a placeholder model
% name; rename globally if you prefer.
% Sections: Problem Formulation, Architecture, Surrogate Calibration,
% Policy Optimization (incl. post-processing local search).
% =========================================================================

\section{Problem Formulation}
\label{sec:problem}

\begin{table*}[t]
\caption{Summary of Notation}
\label{tab:notation}
\centering
\footnotesize
\begin{tabular}{@{}l p{0.295\textwidth} @{\qquad} l p{0.295\textwidth}@{}}
\toprule
Symbol & Meaning & Symbol & Meaning \\
\midrule
\multicolumn{2}{@{}l}{\itshape Problem setting} &
\multicolumn{2}{@{}l}{\itshape Architecture} \\
$\mathcal{W}$ & grid workspace (four-connected unit cells) &
  $\bm{x}^{A}_k,\bm{x}^{J}_j$ & AMR / job input features ($\mathbb{R}^{8},\mathbb{R}^{16}$) \\
$\Delta(u,v)$ & Manhattan distance between positions $u,v$ &
  $h$ & hidden width of all embeddings ($h{=}128$) \\
$\mathcal{J},\ n$ & job set; number of jobs &
  $\bm{h}^{A}_k$ & AMR embedding \\
$\mathcal{M},\ m$ & AMR fleet; fleet size &
  $\bm{g}^{(l)}_j$ & job embedding after $l$ GIN layers \\
$\mathcal{T}$ & set of material types &
  $\mathcal{G}_t=(\mathcal{J},E_t)$ & dynamic precedence graph at step $t$ \\
$\mathcal{D}^{\mathrm{in}},\mathcal{D}^{\mathrm{out}},\mathcal{D}$ & inbound docks; outbound docks; their union &
  $\mathcal{P}(j)$ & committed predecessors of job $j$ in $\mathcal{G}_t$ \\
$\tau_j$ & material type of job $j$ &
  $\epsilon^{(l)},\ L$ & learnable GIN self-weight; number of GIN layers \\
$o_j,\ d_j$ & inbound and outbound dock of job $j$ &
  $\mathrm{LN}$ & layer normalization \\
$r_j$ & release time of job $j$ &
  $\bm{e}_{\omega}$ & learned operation-type embedding \\
$p_j=p(\tau_j)$ & dock service duration of job $j$ &
  $z_{\omega,k,j}$ & logit of action $(\omega,k,j)$ \\
$\pi_j,\ \delta_j$ & pickup and delivery operations of job $j$ &
  $V_{\phi}$ & critic (state-value) network, parameters $\phi$ \\
$Q$ & per-type AMR carrying capacity &
  & \\
$\bm{a},\ a_j$ & AMR assignment vector; AMR assigned to job $j$ &
  \multicolumn{2}{@{}l}{\itshape Training and refinement} \\
$\sigma,\ \sigma_t$ & operation sequence; prefix committed up to step $t$ &
  $R,\ R(\sigma_{1:t})$ & dispatching rule; completion of prefix $\sigma_{1:t}$ by $R$ \\
$s^{\pi}_j,c^{\pi}_j$ & start / completion of pickup service of job $j$ &
  $C[\cdot]$ & shorthand for $C_{\max}(\Phi(\cdot))$ \\
$s^{\delta}_j,c^{\delta}_j$ & start / completion of delivery service of job $j$ &
  $A_t,\ \hat{A}_t$ & stepwise counterfactual advantage; standardized \\
$t_k(u,v)$ & realized travel time of AMR $k$ from $u$ to $v$ &
  $A^{\mathrm{lb}}_t$ & load-balance shaping advantage \\
$\Phi$ & collision-aware executor &
  $\lambda_{\mathrm{lb}},\ \lambda_{H}$ & shaping and entropy coefficients \\
$C_{\max},\ F_k$ & makespan; finish time of AMR $k$ incl.\ depot return &
  $\mathcal{H}_t$ & entropy of the masked policy at step $t$ \\
$W_d$ & number of physical waiting cells at dock $d$ &
  $\mathcal{L}(\theta)$ & training loss \\
\multicolumn{2}{@{}l}{\itshape Sequential decision process} &
  $B,\ E$ & batch size (episodes); number of training epochs \\
$s_t,\ a_t$ & state and action at decision epoch $t$ &
  $T_{\mathrm{val}},\ \kappa$ & validation interval; infeasibility penalty \\
$\omega$ & operation type, $\omega\in\{\textsc{pick},\textsc{deliver}\}$ &
  $I_{\mathrm{cf}},\ I_{\mathrm{ca}}$ & collision-free / collision-aware refinement budgets \\
$\Lambda(s_t)$ & feasible-action set (mask) in state $s_t$ &
  $\mathbf{1}[\cdot]$ & indicator function \\
$\hat{\Psi}$ & calibrated surrogate transition model &
  & \\
$\mathcal{E}_d$ & committed service events $(s_e,c_e,j_e,\omega_e)$ at dock $d$ &
  & \\
$\hat{b}_k,\ x_k$ & next-available time and position of AMR $k$ &
  & \\
$\rho$ & ready time of an AMR at a dock &
  & \\
$\hat{t}(u,v)$ & calibrated travel-time estimate \eqref{eq:travelcal} &
  & \\
$\hat{w}_d(\rho)$ & queue-aware dock-wait estimate \eqref{eq:waitcal} &
  & \\
$q_d(\rho)$ & services still unfinished at time $\rho$ at dock $d$ &
  & \\
$\alpha,\ \beta$ & travel-calibration scale and offset &
  & \\
$\varphi(\cdot)$ & dock-wait penalty table over queue depths &
  & \\
$\pi_{\theta}$ & scheduling policy with parameters $\theta$ &
  & \\
\bottomrule
\end{tabular}
\end{table*}

\subsection{System Description}

We consider the internal transfer problem of a cross-docking terminal
served by a fleet of autonomous mobile robots (AMRs). Table~\ref{tab:notation}
summarizes the notation used throughout the paper. The terminal floor
is modeled as a four-connected unit grid $\mathcal{W}$ in which a set of
\emph{inbound docks} $\mathcal{D}^{\mathrm{in}}$ is located along one
boundary and a set of \emph{outbound docks} (processing stations)
$\mathcal{D}^{\mathrm{out}}$ along the opposite boundary; the AMR depots
lie between them. Consistent with cross-docking practice, the floor is
free of static obstacles, so congestion arises from vehicle interference
and dock contention rather than from the layout itself. AMRs travel at
unit speed between adjacent cells, and the distance between two
positions $u,v$ is the Manhattan metric $\Delta(u,v)$.

A set of transfer jobs $\mathcal{J}=\{1,\dots,n\}$ must be moved across
the terminal by a fleet $\mathcal{M}=\{1,\dots,m\}$ of homogeneous AMRs.
Each job $j\in\mathcal{J}$ represents one unit load of a material type
$\tau_j\in\mathcal{T}$ that becomes available at its inbound dock
$o_j\in\mathcal{D}^{\mathrm{in}}$ at release time $r_j\ge 0$ and must be
delivered to a designated outbound dock
$d_j\in\mathcal{D}^{\mathrm{out}}$. Handling a unit load of type $\tau$
occupies a dock for a type-dependent service duration $p(\tau)$; we
write $p_j=p(\tau_j)$ and assume the loading service at the inbound dock
and the unloading/processing service at the outbound dock take the same
time, reflecting symmetric handling of a unit load. Each job therefore
consists of an ordered pair of operations: a \emph{pickup} operation
$\pi_j$ executed at $o_j$ and a \emph{delivery} operation $\delta_j$
executed at $d_j$, both performed by the same AMR (no transshipment).
An AMR may consolidate loads subject to a per-type carrying capacity
$Q$: at no time may it carry more than $Q$ units of any single material
type.

Docks are exclusive single-server resources: at most one AMR may be in
service at a dock at any time, and services are non-preemptive. An AMR
arriving at an occupied dock queues in a bounded set of waiting cells
adjacent to that dock; when the waiting line is saturated, the AMR must
hold at an upstream position and re-approach later. Finally, AMR motion
must be collision-free: at any unit time step, each grid cell may be
occupied by at most one AMR.

\subsection{Formal Statement}

A solution is a pair $(\bm{a},\sigma)$, where
$\bm{a}=(a_1,\dots,a_n)\in\mathcal{M}^{n}$ assigns each job to an AMR
and $\sigma$ is a total order over the $2n$ operations
$\{\pi_j,\delta_j\}_{j\in\mathcal{J}}$ that determines the sequence in
which operations are committed. Let $s^{\pi}_j,c^{\pi}_j$
(resp.\ $s^{\delta}_j,c^{\delta}_j$) denote the service start and
completion times of $\pi_j$ (resp.\ $\delta_j$) induced by executing
$(\bm{a},\sigma)$. A feasible execution satisfies:
\begin{align}
c^{\pi}_j &= s^{\pi}_j + p_j, \qquad
c^{\delta}_j = s^{\delta}_j + p_j,
&& \forall j\in\mathcal{J}, \label{eq:nonpreempt}\\
s^{\pi}_j &\ge r_j,
&& \forall j\in\mathcal{J}, \label{eq:release}\\
s^{\delta}_j &\ge c^{\pi}_j,
&& \forall j\in\mathcal{J}, \label{eq:precedence}
\end{align}
\begin{align}
\big|\{\,j : a_j{=}k,\ \tau_j{=}\tau,\ c^{\pi}_j \le u < c^{\delta}_j \}\big| &\le Q,
&& \forall k,\tau,u, \label{eq:capacity}\\
\big[s_e, c_e\big) \cap \big[s_{e'}, c_{e'}\big) &= \emptyset,
&& \forall e\ne e' \ \text{at the same dock}, \label{eq:exclusive}
\end{align}
where \eqref{eq:nonpreempt} enforces non-preemptive services,
\eqref{eq:release} release dates, \eqref{eq:precedence} pickup-before-%
delivery precedence, \eqref{eq:capacity} the per-type carrying capacity,
and \eqref{eq:exclusive} dock exclusivity over all service events $e$.
In addition, for any two operations $u\prec v$ executed consecutively by
the same AMR $k$,
\begin{equation}
s_v \;\ge\; c_u + t_k\big(\mathrm{loc}(u),\mathrm{loc}(v)\big),
\label{eq:travel}
\end{equation}
where $t_k(\cdot,\cdot)\ge\Delta(\cdot,\cdot)$ is the realized travel
time of AMR $k$, which may exceed the Manhattan distance due to
collision avoidance, queuing detours, and dock-clearing maneuvers.

The realized travel and waiting times are not fixed constants but are
determined by a high-fidelity \emph{executor} $\Phi$, which routes every
AMR leg with space--time $\mathrm{A}^{*}$ search over a shared
reservation table (at most one AMR per cell per time step), manages the
bounded dock waiting lines, and inserts upstream holding maneuvers when
lines are saturated. The problem is therefore a simulation-based
optimization: find
\begin{equation}
(\bm{a}^{*},\sigma^{*}) \;=\;
\arg\min_{(\bm{a},\sigma)} \; C_{\max}\big(\Phi(\bm{a},\sigma)\big),
\label{eq:objective}
\end{equation}
where the makespan $C_{\max}=\max_{k\in\mathcal{M}} F_k$ and $F_k$ is
the time at which AMR $k$ completes its final activity, including its
return trip to the depot. The problem generalizes both the capacitated
pickup-and-delivery problem and parallel-machine scheduling with
sequence-dependent setups, and is therefore NP-hard; the coupling to
collision-free multi-agent routing further compounds its difficulty
\cite{stern2019mapf}.

\subsection{Sequential Decision Formulation}
\label{sec:mdp}

We cast the construction of $(\bm{a},\sigma)$ as an episodic sequential
decision process with exactly $2n$ decision epochs, one per committed
operation. At epoch $t$, the state
\begin{equation}
s_t=\big(\{\bm{x}^{A}_{k}\},\, \{\bm{x}^{J}_{j}\},\,
\{\mathcal{E}_d\},\, \sigma_{t}\big)
\label{eq:state}
\end{equation}
collects the AMR states (position, next-available time, per-type
inventory), the job states (unpicked / onboard / delivered, carrier
identity), the set of committed dock service events
$\mathcal{E}_d=\{(s_e,c_e,j_e,\omega_e)\}$ at each dock $d$, and the
partial operation sequence $\sigma_t$.

An action is a triple
\begin{equation}
a_t=(\omega,k,j)\in\{\textsc{pick},\textsc{deliver}\}\times\mathcal{M}\times\mathcal{J},
\label{eq:action}
\end{equation}
giving a flat action space of size $2mn$ restricted by a feasibility
mask $\Lambda(s_t)$: $(\textsc{pick},k,j)$ is legal iff $j$ is unpicked
and AMR $k$ carries fewer than $Q$ units of type $\tau_j$;
$(\textsc{deliver},k,j)$ is legal iff $j$ is onboard AMR $k$. The mask
guarantees by construction that every complete episode yields a feasible
solution, so no repair mechanism is required.

Transitions follow a deterministic surrogate model $\hat{\Psi}$ that
advances the acting AMR's clock using a \emph{calibrated} travel-time
estimate $\hat{t}(\cdot,\cdot)$ and a queue-aware dock-wait estimate
$\hat{w}_d(\cdot)$ fitted offline against the executor $\Phi$
(Section~\ref{sec:calibration}); committing an action appends the
corresponding service window to $\mathcal{E}_d$. For a pickup, the
service start of job $j$ by AMR $k$ located at $x_k$ with next-available
time $\hat{b}_k$ is
\begin{equation}
s^{\pi}_j=\underbrace{\max\!\big(\hat{b}_k+\hat{t}(x_k,o_j),\, r_j\big)}_{\text{ready time }\rho}
\;+\;\hat{w}_{o_j}(\rho),
\label{eq:surrogate}
\end{equation}
and analogously for deliveries at $d_j$. After $2n$ steps the completed
solution is evaluated by the executor, and the learning objective is to
minimize the expected executed makespan
$\mathbb{E}_{\pi_\theta}\!\left[C_{\max}(\Phi(\bm{a},\sigma))\right]$;
the policy-gradient training procedure and its dispatch-rule baseline
are described in Section~\ref{sec:training}.

% =========================================================================

\section{Sparse-Graph Multi-Pointer Dispatcher}
\label{sec:architecture}

\subsection{Overview and Design Rationale}

Learning-to-dispatch architectures for shop scheduling encode
operations as nodes of a disjunctive graph and learn a construction
policy over it \cite{zhang2020l2d,song2023fjsp}. Transferring this
paradigm to AMR-served cross-docking raises a representation question:
how should the tight resource coupling---$m$ vehicles contending for
$|\mathcal{D}|$ single-server docks under collision-constrained
motion---enter the graph? A seemingly natural answer is to join every
pair of jobs that share a dock. We found this to be
counterproductive: with $n/|\mathcal{D}|$ jobs per dock, such
contention edges form dense cliques, and message passing over them
drives all job embeddings toward a common average (pairwise cosine
similarity above $0.99$ at initialization in our measurements),
leaving the policy unable to discriminate between candidate jobs
precisely at the opening decisions that dominate the makespan.
Over-smoothing under dense neighborhoods is a known failure mode of
deep graph networks \cite{li2018deeper,oono2020oversmoothing}; in this
problem it appears already at three message-passing rounds.

The proposed scheduler, hereafter \textsc{CD-GNN}, therefore keeps the
graph \emph{sparse and dynamic}: edges encode only scheduling
precedences that have actually been committed, while resource
congestion enters through calibrated, decision-grade node features.
The model comprises (i)~a job encoder over a dynamic
\emph{precedence graph} processed by residual graph isomorphism
network (GIN) layers \cite{xu2019gin}; (ii)~a lightweight fleet
encoder; and (iii)~a factored multi-pointer actor that scores each
masked (operation, AMR, job) triple \cite{kool2019attention}. The full
model is deliberately lean: $2.8\times10^{5}$ parameters
($2.5\times10^{5}$ actor, $3.3\times10^{4}$ critic) with hidden width
$h=128$ throughout, permutation-invariant in the job and AMR sets, and
applicable unchanged to any instance size.

% TODO: architecture overview figure
% \begin{figure*}[t]
%   \centering
%   \includegraphics[width=\textwidth]{figures/cdgnn_architecture}
%   \caption{Overview of \textsc{CD-GNN}. Job features (including
%   calibrated congestion estimates) are encoded and refined over the
%   dynamic precedence graph by residual GIN layers; AMR features are
%   encoded independently; the multi-pointer actor scores each feasible
%   (operation, AMR, job) triple under the feasibility mask.}
%   \label{fig:architecture}
% \end{figure*}

\subsection{Congestion-Aware State Representation}
\label{sec:features}

At every decision epoch the state \eqref{eq:state} is rendered into two
feature families, summarized in Table~\ref{tab:features}. All features
are normalized to comparable magnitudes by fixed scales.

\emph{Job nodes} ($\bm{x}^{J}_j\in\mathbb{R}^{16}$). Each job encodes
its service duration, lifecycle status (unpicked/onboard/delivered) and
carrier identity, its material type, and the coordinates of both its
supply position and its destination station, giving the network direct
geometric grounding. Resource congestion enters through four derived
features: the current availability delays of the job's inbound and
outbound docks, and the number of still-unpicked competitor jobs
committed to each of the two docks. A calibrated completion-time
estimate (Section~\ref{sec:calibration}) summarizes how quickly the job
could finish if served next, providing a value-like signal that the
network would otherwise need several layers to reconstruct.

\emph{AMR nodes} ($\bm{x}^{A}_k\in\mathbb{R}^{8}$). Each AMR encodes a
busy indicator and its time until available, per-type inventory
ratios, its coordinates, and the number of jobs currently assigned to
it, which exposes fleet-utilization imbalance to the policy.

\begin{table}[t]
\caption{Node Features of \textsc{CD-GNN}}
\label{tab:features}
\centering
\footnotesize
\begin{tabular}{@{}llc@{}}
\toprule
& Feature & Scale \\
\midrule
\multirow{12}{*}{\rotatebox{90}{Job $\bm{x}^{J}_j\!\in\!\mathbb{R}^{16}$}}
& existence flag & -- \\
& service duration $p_j$ & $1/25$ \\
& carrier index (0 if unpicked) & $1/m$ \\
& destination position $(x,y)$ ($\times2$) & $1/10$ \\
& supply position $(x,y)$ ($\times2$) & $1/10$ \\
& material type one-hot ($\times3$) & -- \\
& lifecycle status $\{0,\tfrac12,1\}$ & -- \\
& completion-time estimate & $1/500$ \\
& inbound dock availability delay & $1/100$ \\
& outbound dock availability delay & $1/100$ \\
& inbound dock committed queue & $1/m$ \\
& outbound dock committed queue & $1/m$ \\
\midrule
\multirow{7}{*}{\rotatebox{90}{AMR $\bm{x}^{A}_k\!\in\!\mathbb{R}^{8}$}}
& busy indicator & -- \\
& time until available & $1/50$ \\
& inventory ratio, per type ($\times3$) & $1/Q$ \\
& position $(x,y)$ ($\times2$) & $1/10$ \\
& assigned-job count & $1/10$ \\
\bottomrule
\end{tabular}
\end{table}

\subsection{Dynamic Precedence Graph and Job Encoder}
\label{sec:graph}

Job features pass through a linear map
$\bm{g}^{(0)}_j=\bm{W}_{\!J}\,\bm{x}^{J}_j\in\mathbb{R}^{h}$ and AMR
features through a linear encoder
$\bm{h}^{A}_k=\bm{W}_{\!A}\,\bm{x}^{A}_k\in\mathbb{R}^{h}$.

Jobs are coupled by the \emph{committed} structure of the partial
schedule. Following the adding-arc construction of disjunctive-graph
dispatching \cite{zhang2020l2d}, the graph
$\mathcal{G}_t=(\mathcal{J},E_t)$ starts empty and grows one or two
directed arcs per decision: when an operation of job $j$ is committed,
an arc is added from the previously scheduled job on the same AMR to
$j$, and from the previously scheduled job on the same outbound dock to
$j$. Each job thus aggregates information from its immediate
predecessors $\mathcal{P}(j)$ on both resources it occupies---the
vehicle sequence and the dock sequence---while uncommitted jobs
interact only through the shared congestion features of
Section~\ref{sec:features}, never through dense same-dock edges.

Message passing applies $L=3$ residual GIN layers with layer
normalization,
\begin{equation}
\bm{g}^{(l+1)}_j=\bm{g}^{(l)}_j+
\mathrm{LN}\!\left(\mathrm{MLP}^{(l)}\!\Big(
(1{+}\epsilon^{(l)})\,\bm{g}^{(l)}_j+
\sum_{i\in\mathcal{P}(j)}\!\bm{g}^{(l)}_i
\Big)\right)\!,
\label{eq:gin}
\end{equation}
retaining the injective sum aggregation of GIN \cite{xu2019gin}: with
at most two predecessors per node, message magnitudes stay bounded
irrespective of instance size, so the scale-invariance concerns that
arise with dense contention edges do not apply.

\subsection{Multi-Pointer Action Scoring and Value Estimation}
\label{sec:actor}

Let $\bm{e}_{\omega}\in\mathbb{R}^{h}$ be a learned embedding of the
operation type. The actor scores every candidate triple by
\begin{equation}
z_{\omega,k,j}=\mathrm{MLP}_{\pi}\!\big(\big[\,\bm{e}_{\omega}\,\|\,
\bm{g}^{(L)}_{j}\,\|\,\bm{h}^{A}_{k}\big]\big),
\label{eq:logit}
\end{equation}
and the policy is the masked softmax
\begin{equation}
\pi_{\theta}(a\,|\,s_t)=
\frac{\exp(z_{a})\,\mathbf{1}[a\in\Lambda(s_t)]}
{\sum_{a'\in\Lambda(s_t)}\exp(z_{a'})}.
\label{eq:policy}
\end{equation}
The critic estimates the state value from pooled entity summaries,
\begin{equation}
V_{\phi}(s_t)=\mathrm{MLP}_{V}\!\big(\big[\,
\overline{\bm{g}}^{\mathrm{(unf)}} \,\|\, \overline{\bm{h}}^{A}\big]\big),
\label{eq:critic}
\end{equation}
where $\overline{\bm{g}}^{\mathrm{(unf)}}$ averages the embeddings of
unfinished jobs only. A full episode requires exactly $2n$ forward
passes; each pass scores all $2mn$ actions in a single batched
evaluation of \eqref{eq:logit}.

% =========================================================================

\section{Surrogate Model Calibration}
\label{sec:calibration}

\subsection{The Two-Simulator Gap}

The policy constructs schedules on the surrogate $\hat{\Psi}$ of
Section~\ref{sec:mdp}, whereas solution quality is defined by the
executor $\Phi$ of \eqref{eq:objective}. An uncalibrated surrogate that
advances clocks by Manhattan distances and single-scalar dock
availabilities is systematically optimistic: it ignores collision
detours, dock-clearing maneuvers, waiting-line approach paths, and the
upstream holds triggered when waiting lines saturate. The congestion
features of Section~\ref{sec:features}---and therefore every decision
the policy makes---would then be conditioned on clocks that drift
further from reality as congestion grows, precisely in the regime where
scheduling decisions matter most. Rather than paying the computational
cost of embedding the full executor in the training loop, we close most
of this gap with two lightweight corrections fitted offline.

\subsection{Calibrated Travel and Dock-Wait Estimators}

\emph{Travel.} Realized leg times are modeled as an affine function of
the Manhattan distance, clamped from below by the metric itself:
\begin{equation}
\hat{t}(u,v)=\max\!\big(\Delta(u,v),\ \alpha\,\Delta(u,v)+\beta\big).
\label{eq:travelcal}
\end{equation}

\emph{Dock waiting.} The wait experienced by an AMR that becomes ready
at a dock $d$ at time $\rho$ is modeled as the committed availability
delay plus a penalty indexed by the queue depth at readiness,
\begin{equation}
\hat{w}_d(\rho)=\max\!\big(0,\ \mathrm{avail}_d-\rho\big)
+\varphi\!\big(\min(q_d(\rho),\,3)\big),
\label{eq:waitcal}
\end{equation}
where $q_d(\rho)=|\{e\in\mathcal{E}_d : c_e>\rho\}|$ counts committed
services unfinished at $\rho$, and
$\varphi:\{0,1,2,3^{+}\}\to\mathbb{R}_{\ge0}$ is a monotone lookup
table. The bounded table shape is deliberate: waiting-line detours and
upstream holds produce a cost cliff once the line saturates, a
nonlinearity that a linear congestion model cannot express.

\subsection{Fitting Procedure}

Calibration data are obtained by replaying complete schedules through
both simulators and aligning them per operation. Schedules are drawn
from heterogeneous sources---five dispatching-rule pairs and the
current learned policies---across instance sizes
$n\in\{10,20,40,60\}$, so that the fitted correction covers the state
distribution that policies actually visit. For every operation, the
raw surrogate supplies the Manhattan leg distance, the base wait, and
the queue depth at readiness; the executor's logged timeline is
partitioned, per AMR, into legs terminating at each service start,
within which movement entries sum to the realized travel and the
remainder is attributed to waiting. Executions containing infeasible
operations are discarded. The travel parameters $(\alpha,\beta)$ are
fitted by least squares and the penalties $\varphi(\cdot)$ as clipped
bin means with monotonicity enforced across depths.

In our environment, fitting on $10{,}360$ operation pairs from $159$
schedules yielded $\alpha=0.963$, $\beta=1.12$ ($R^{2}=0.91$) and
$\varphi=(0.04,\,0.04,\,1.66,\,3.57)$; the step between depths one and
two locates the waiting-line saturation cliff. The correction removes
the mean per-leg travel bias entirely ($+0.744\to0.000$ time units)
and halves the end-to-end makespan drift between surrogate and
executor (mean absolute error $24.3\to12.1$ on held-out schedules).
The calibration is a frozen artifact: it is fitted once before
training, adds no learned components to the training loop, and is
shared by \emph{all} policies and by the dispatching-rule baseline of
Section~\ref{sec:training}, so every method plans in the same
corrected world. After training converges, the fit can be re-estimated
from the new policies' schedules to confirm stability.

% =========================================================================

\section{Policy Optimization and Refinement}
\label{sec:training}

\subsection{Objective and Gradient Estimator}

The policy $\pi_\theta$ is trained to minimize the expected executed
makespan $\mathbb{E}_{\pi_\theta}[C_{\max}(\Phi(\bm{a},\sigma))]$ with
the REINFORCE estimator \cite{williams1992reinforce}. During training,
episodes are generated by sampling from the masked policy
\eqref{eq:policy} on the calibrated surrogate; at validation and
deployment, actions are selected greedily. The gradient quality of
REINFORCE hinges on its baseline; we use a counterfactual baseline
built from a strong dispatching rule and evaluated by the executor
itself.

\subsection{Counterfactual Dispatch-Rule Baseline}
\label{sec:baseline}

Let $R$ denote a deterministic dispatching heuristic (job rule and AMR
rule; we use earliest-completion for both) and let
$R(\sigma_{1:t})$ denote the complete solution obtained by executing
the partial sequence $\sigma_{1:t}$ and completing all remaining
operations with $R$. Define the shorthand
$C[\cdot]=C_{\max}(\Phi(\cdot))$. The advantage of the action taken at
step $t$ is the improvement of committing that action and then
following the rule, relative to following the rule immediately:
\begin{equation}
A_t \;=\; C\big[R(\sigma_{1:t-1})\big]\;-\;
C\big[R(\sigma_{1:t-1}\oplus a_t)\big].
\label{eq:advantage}
\end{equation}
Both completions are evaluated by the collision-aware executor, so the
learning signal is grounded in true solution quality rather than in
surrogate clocks. Unlike episode-level rollout baselines
\cite{rennie2017self,kool2019attention}, \eqref{eq:advantage} assigns
credit per decision: an early pickup that congests a dock is penalized
at the step where it is committed, not diluted across the episode.
The price is $\mathcal{O}(n)$ executor evaluations per episode, which
dominates training cost; we consider the sharper credit assignment a
favorable trade at the instance sizes studied here, and
Section~\ref{sec:training-details} lists a cheaper actor--critic
alternative.

\subsection{Loss Function}

Advantages are standardized to zero mean and unit variance across all
steps of a batch of $B$ episodes. The per-batch loss combines the
policy gradient with a load-balancing shaping term and an entropy
bonus:
\begin{equation}
\mathcal{L}(\theta)=
-\frac{1}{B}\sum_{b=1}^{B}\sum_{t=1}^{2n}
\Big[\log\pi_\theta(a_t^b|s_t^b)\,
\big(\hat{A}_t^{\,b}+\lambda_{\mathrm{lb}}A_t^{\mathrm{lb},b}\big)
+\lambda_{H}\,\mathcal{H}_t^{b}\Big],
\label{eq:loss}
\end{equation}
where $\hat{A}_t$ is the standardized advantage,
$\mathcal{H}_t$ is the entropy of the masked policy over feasible
actions, and the shaping term
$A^{\mathrm{lb}}_t=\min_{k'}c_{k'}-c_{k}$ (for pickups; zero for
deliveries) compares the assignment count $c_k$ of the chosen AMR with
the least-loaded AMR, discouraging degenerate fleet utilization early
in training \cite{ng1999shaping}. We use
$\lambda_{\mathrm{lb}}=0.1$ and $\lambda_{H}=0.01$.

\begin{algorithm}[t]
\caption{Training with counterfactual dispatch baseline}
\label{alg:training}
\begin{algorithmic}[1]
\Require policy $\pi_\theta$, rule $R$, executor $\Phi$, surrogate $\hat{\Psi}$
\State fit calibration $(\alpha,\beta,\varphi)$ offline; freeze (Sec.~\ref{sec:calibration})
\For{$\mathrm{epoch}=1,\dots,E$}
  \For{$b=1,\dots,B$}
    \State sample instance; roll out $\sigma^b\!\sim\!\pi_\theta$ on $\hat{\Psi}$
    \For{$t=1,\dots,2n$}
      \State $A^b_t \gets C[R(\sigma^b_{1:t-1})]-C[R(\sigma^b_{1:t-1}\oplus a^b_t)]$
    \EndFor
  \EndFor
  \State standardize $\{A^b_t\}$ over the batch
  \State update $\theta$ by Adam on $\mathcal{L}(\theta)$ in \eqref{eq:loss}; clip $\lVert\nabla\rVert\le1$
  \If{$\mathrm{epoch}\bmod T_{\mathrm{val}}=0$}
    \State greedy-decode held-out instances; score by
    $C_{\max}+\kappa\cdot(\text{infeasible ops})$; keep best
  \EndIf
\EndFor
\end{algorithmic}
\end{algorithm}

\subsection{Training Protocol}
\label{sec:training-details}

Algorithm~\ref{alg:training} summarizes the procedure. Only the actor
parameters (encoders, GIN layers, operation embedding, and actor head)
receive gradients under REINFORCE; the critic \eqref{eq:critic} is
reserved for an optional proximal-policy-optimization variant
\cite{schulman2017ppo} that replaces \eqref{eq:advantage} with
critic-based advantages when the executor evaluations of the stepwise
baseline are prohibitive. We train with Adam \cite{kingma2015adam}
(learning rate $3\times10^{-4}$, batch $B=16$, gradient-norm clip
$1.0$) for $2{,}000$ epochs. Every $T_{\mathrm{val}}=50$ epochs the
policy is evaluated by greedy decoding on a fixed held-out validation
set, disjoint from the training pool, and scored by the executed
makespan plus a penalty of $\kappa=10^{3}$ per infeasible operation;
the checkpoint with the best validation score is retained. Computing
the entropy in \eqref{eq:loss} over the feasible-action support only
keeps the bonus comparable as the mask tightens toward the end of an
episode.

\subsection{Post-Processing Local Search}
\label{sec:localsearch}

At deployment, the greedy-decoded schedule is refined by a two-stage
first-improvement local search that reuses the neighborhood operators
of our genetic-algorithm baseline, so that quality differences between
methods are attributable to the constructed solutions rather than to
the refinement. A move either swaps two operations in $\sigma$, at
least one of which belongs to the AMR that currently determines the
makespan, or relocates a contiguous block of same-AMR, same-type
operations; the AMR assignment $\bm{a}$ is left unchanged. The first
stage performs $I_{\mathrm{cf}}=10^{3}$ trials under the fast
collision-free evaluator to explore broadly; the second performs
$I_{\mathrm{ca}}=10^{2}$ trials under the full executor $\Phi$ to
polish under true routing. The refinement adds approximately two
seconds per instance and is entirely optional: all learned-policy
results are reported both with and without it.

% =========================================================================
% Suggested BibTeX entries (replace keys above with your bibliography):
%
% @inproceedings{xu2019gin, Xu, Hu, Leskovec, Jegelka, "How Powerful Are
%   Graph Neural Networks?", ICLR 2019}
% @inproceedings{zhang2020l2d, Zhang et al., "Learning to Dispatch for
%   Job Shop Scheduling via Deep Reinforcement Learning", NeurIPS 2020}
% @article{song2023fjsp, Song et al., "Flexible Job-Shop Scheduling via
%   Graph Neural Network and Deep Reinforcement Learning", IEEE TII 2023}
% @inproceedings{li2018deeper, Li, Han, Wu, "Deeper Insights into Graph
%   Convolutional Networks for Semi-Supervised Learning", AAAI 2018}
% @inproceedings{oono2020oversmoothing, Oono, Suzuki, "Graph Neural
%   Networks Exponentially Lose Expressive Power for Node
%   Classification", ICLR 2020}
% @inproceedings{stern2019mapf, Stern et al., "Multi-Agent Pathfinding:
%   Definitions, Variants, and Benchmarks", SoCS 2019}
% @article{williams1992reinforce, Williams, "Simple Statistical
%   Gradient-Following Algorithms for Connectionist Reinforcement
%   Learning", Machine Learning 8, 1992}
% @inproceedings{rennie2017self, Rennie et al., "Self-Critical Sequence
%   Training for Image Captioning", CVPR 2017}
% @inproceedings{kool2019attention, Kool, van Hoof, Welling, "Attention,
%   Learn to Solve Routing Problems!", ICLR 2019}
% @inproceedings{ng1999shaping, Ng, Harada, Russell, "Policy Invariance
%   Under Reward Transformations: Theory and Application to Reward
%   Shaping", ICML 1999}
% @article{schulman2017ppo, Schulman et al., "Proximal Policy
%   Optimization Algorithms", arXiv:1707.06347}
% @inproceedings{kingma2015adam, Kingma, Ba, "Adam: A Method for
%   Stochastic Optimization", ICLR 2015}
% =========================================================================
